\documentclass[11pt]{article}
\pdfinfoomitdate=1
\pdftrailerid{}
\pdfsuppressptexinfo=-1
\usepackage[a4paper,margin=28mm]{geometry}
\usepackage[T1]{fontenc}
\usepackage{lmodern,amsmath,amssymb,amsthm,mathtools,microtype}
\usepackage{xurl,hyperref,enumitem}
\hypersetup{colorlinks=true,linkcolor=blue,citecolor=blue,urlcolor=blue,
 pdftitle={Dimension structure in simplex stability},
 pdfauthor={Entry 005 AI-assisted research manuscript candidate}}
\setlength{\emergencystretch}{2em}
\numberwithin{equation}{section}
\newtheorem{theorem}{Theorem}[section]
\newtheorem{lemma}[theorem]{Lemma}
\newtheorem{proposition}[theorem]{Proposition}
\newtheorem{corollary}[theorem]{Corollary}
\theoremstyle{definition}\newtheorem{definition}[theorem]{Definition}
\theoremstyle{remark}\newtheorem{remark}[theorem]{Remark}
\newcommand{\R}{\mathbb R}
\newcommand{\E}{\mathbb E}
\newcommand{\conv}{\operatorname{conv}}
\newcommand{\vol}{\operatorname{Vol}}
\newcommand{\pyr}{\operatorname{pyr}}
\newcommand{\supp}{\operatorname{supp}}
\newcommand{\diag}{\operatorname{diag}}
\newcommand{\Id}{I}
\newcommand{\ones}{\mathbf 1}
\newcommand{\ip}[2]{\langle #1,#2\rangle}
\newcommand{\norm}[1]{\lVert #1\rVert}
\newcommand{\abs}[1]{\lvert #1\rvert}
\title{Dimension structure in simplex stability}
\author{Entry 005 research manuscript\\\small AI-assisted manuscript candidate}
\date{8 October 2026}
\begin{document}
\maketitle
\begin{abstract}
We study approximation of a convex body by each prescribed maximum-volume
inscribed simplex, with dilation centered at that simplex's own centroid,
using the projection-body pyramid defect. For polytopes with at most
$d+2$ facets, the optimal coefficient of the sharp power $1/(d-1)$ has
asymptotic value $d$: its ratio to $d$ tends to one. The geometric basis is
an exact classification of all maximum simplices in a simplex cut by one
halfspace, including nonvertex maxima, through weighted rooted forests.
Signed maximal cofactors of the actual facet area-normal/support matrix
give an exact formula for the defect. We also derive concentric cores and
a face-depth hierarchy from full projection control, and prove
linear-dimensional retention for edge-supported configurations through
inverse-minor rigidity of stochastic matrices with two-supported columns.
We reproduce the finite actual-facet enclosure argument and preserve its
compatibility conditions. The unrestricted gap between linear lower and
quadratic upper dimension orders remains open.
\end{abstract}
\noindent\textbf{Status.} This is an unpublished manuscript candidate.
The new results have written proofs and the source-specific traditional
review record described in Appendix~\ref{app:review}; they have not been
formalized in Lean. Independent model review, parent integration review,
and an author's self-reviewed contribution are distinguished there.
No professional peer-review, novelty certification, or priority claim is made.
\setcounter{tocdepth}{1}
\tableofcontents
\section{The original invariant and the dimension question}\label{sec:intro}
For a full-dimensional convex body $K\subset\R^d$, write $|K|$ for its
Euclidean volume and $h_K$ for its support function. For a unit vector $u$,
let $\pi_K(u)$ denote the $(d-1)$-dimensional volume of the orthogonal
projection onto $u^\perp$. Its homogeneous extension is
\[
 b_K(z)=\begin{cases}\norm z\pi_K(z/\norm z),&z\ne0,\\0,&z=0.\end{cases}
\]
The canonical projection body $\Pi K$ has support function $b_K$.
For a body $L$ of dimension $p$, put
\[
 R(L)=\frac{|\Pi L|}{|L|^{p-1}},\qquad
 \pyr L=\conv((L,0)\cup\{(0,1)\}).
\]
The actual projection-body pyramid quantities are
\begin{equation}\label{eq:actualdefect}
 a(K)=\left(\frac d{d+1}\right)^d\frac{R(\pyr K)}{R(K)}-1,
 \qquad e(K)=a(K)-\frac1{d+1}.
\end{equation}
For a full-dimensional simplex $S\subset K$ with its own centroid $g_S$,
we define the excess
\begin{equation}\label{eq:excess}
 E(K,S)=\inf\{u\ge0:K\subset g_S+(1+u)(S-g_S)\}.
\end{equation}
The simplex $S$ is prescribed. In particular, the definition does not
minimize over maximum simplices or over centers. All affine changes below
transport $K$, $S$, and $g_S$ together. Set
\[
 n=d+1,\qquad m=d-1,\qquad H_d=3d^2(d+1)^3(d+2).
\]
Unless stated otherwise $d\ge3$. The geometric clipping results also apply
in smaller dimensions within their explicitly stated ranges.

The retained unrestricted theorem~\cite{Quadratic} gives
$E(K,S)\le4096d^2e(K)^{1/m}$ for every prescribed maximum simplex;
realized square-pyramid bodies give a lower coefficient of
$n[dn]^{1/m}$. Its proof and formal-verification record are prior inputs,
not formal verification of the statements added here. Our principal result
settles the leading dimensional coefficient on an entire facet-count class.

\begin{theorem}[The few-facet coefficient]\label{thm:few-facet}
For every full-dimensional convex $d$-polytope $K$ with at most $d+2$
facets and every maximum-volume inscribed simplex $S$,
\begin{equation}\label{eq:few-upper}
 E(K,S)\le nH_d^{1/m}e(K)^{1/m}\le192d\,e(K)^{1/m}.
\end{equation}
The power $1/m$ is sharp in this class. If
\[
 G_d^{\rm fac}=\sup\left\{\frac{E(K,S)}{e(K)^{1/m}}:
 K\text{ has at most }d+2\text{ facets},\ S\text{ is maximum},\ e(K)>0\right\},
\]
where all bodies are full dimensional, then
\begin{equation}\label{eq:few-squeeze}
 n[dn]^{1/m}\le G_d^{\rm fac}\le nH_d^{1/m},
 \qquad\lim_{d\to\infty}\frac{G_d^{\rm fac}}d=1.
\end{equation}
\end{theorem}
The exact fixed-dimensional optimum is not determined. The proof in
Section~\ref{sec:few} uses two exact descriptions: all maxima in an arbitrary
one-halfspace clipping, and the actual defect from signed facet cofactors.
The halfspace may remove several vertices of its original simplex.

Full projection control has additional content beyond near-corner points.
Section~\ref{sec:projection} develops inscribed concentric cores and a
complete face-depth hierarchy. The resulting comparison simplex is not
identified with a prescribed maximum. Section~\ref{sec:inverse} proves
inverse rigidity and retention when columns have support at most two;
its actual-defect conclusion requires compatibility with the enclosing
simplex supplied by the actual chain. These conclusions leave the
unrestricted linear-versus-quadratic dimension-order question open.

The actual facet identities, finite determinant assignment, and enclosure
proof are reproduced in Appendices~\ref{sec:assignment}--\ref{sec:enclosure}.
Their standard convex-geometric inputs are Cauchy's projection formula,
the zonotope volume formula, the first mixed-volume formula, and
Minkowski's first inequality; Brunn--Minkowski is also used for one of the
core proofs. We use these classical results in their ordinary forms
\cite{Schneider}. A systematic literature and priority comparison has not
been carried out. Source-series provenance~\cite{Series} and the precise review
boundaries are recorded separately in Appendix~\ref{app:review} and the
accompanying source map.

\section{All maximum simplices in a single halfspace clipping}\label{sec:clipping}
\begin{theorem}[Exact clipping retention]\label{thm:clipping}
Let $P$ be a full-dimensional $d$-simplex, let $H$ be a closed halfspace,
and suppose $K=P\cap H$ is full dimensional. With $n=d+1$ and the
$P$-barycentric coordinates $\lambda_i$, put
\[
 \rho_* =\max_i\left(1-\max_{x\in K}\lambda_i(x)\right).
\]
Every prescribed maximum-volume simplex $S\subset K$ satisfies
\begin{equation}\label{eq:clip-bound}
 E(K,S)\le n\rho_*.
\end{equation}
Consequently corner losses at most $\rho\ge0$ imply $E(K,S)\le n\rho$.
The coefficient $n$ is sharp for every $d\ge2$ and every $0<\rho_*<1$.
All maximum simplices and their exact excess are described below.
\end{theorem}

\subsection{The vertices and a candidate maximum}
Use affine $P$-coordinates, so $P=\conv(e_1,\ldots,e_n)$ in
$\sum_i x_i=1$. Write $H=\{f\ge0\}$ and $f_i=f(e_i)$. If no $f_i$ is
negative, $K=P$ and the result is immediate. Otherwise full dimension
implies that at least one $f_i$ is positive. Partition the indices as
\[
 I=\{i:f_i=-a_i<0\},\qquad J=\{j:f_j=b_j>0\},\qquad Z=\{z:f_z=0\}.
\]
Let $B=\max_{j\in J}b_j>0$, $J_* =\{j:b_j=B\}$, and
\[
 t_i=\frac B{a_i+B},\qquad s_i=\frac{a_i}{a_i+B}=1-t_i\quad(i\in I).
\]
The vertices of $K$ are precisely
\begin{equation}\label{eq:clip-vertices}
 e_j\ (j\in J\cup Z),\qquad
 v_{ij}=\frac{b_je_i+a_ie_j}{a_i+b_j}\quad(i\in I,\ j\in J).
\end{equation}
Indeed, a new vertex lies on an old edge crossing the cut hyperplane.
Equivalently, the active coordinate constraints leave at most two positive
coordinates, and the two-coordinate case gives this expression.

For $z\in Z$, only the vertex $e_z$ has a nonzero $z$-coordinate, so it
is mandatory in every full-dimensional vertex simplex. Its column can be
expanded out. The determinant calculation may therefore be performed on
$I\cup J$ and all $e_z$ then adjoined. Fix $j_*\in J_*$. The simplex
with every retained $e_j$ and with $v_{i j_*}$ for every $i\in I$ has
normalized determinant
\begin{equation}\label{eq:clip-D}
 D_* =\prod_{i\in I}t_i>0.
\end{equation}
Expansion in the retained standard columns proves this identity. Moreover,
\begin{equation}\label{eq:clip-loss}
 \max_K\lambda_i=t_i\ (i\in I),\qquad
 \max_K\lambda_j=1\ (j\in J\cup Z),\qquad
 \rho_* =\max_{i\in I}s_i.
\end{equation}

\subsection{Weighted rooted forests and equality}
Select $|I|+|J|$ columns from the vertices with these row indices.
Represent each $v_{ij}$ by an edge $ij$ in the bipartite graph on $I\cup J$
and each selected $e_j$ by a root mark at $j$. Include isolated row
vertices. Multiply row $i\in I$ by $a_i$ and row $j\in J$ by $-b_j$.
The edge and root columns become, respectively,
\[
 \frac{a_ib_j}{a_i+b_j}(e_i-e_j),\qquad -b_je_j.
\]
Columns in a connected component form a basis of its coordinate space
exactly when that component is a tree with one selected root. An unrooted
component on $v$ vertices has incidence rank $v-1$. A rooted component
has rank at most $v$, and full independence together with the total column
count requires exactly $v$ columns in each component. Connectivity needs
at least $v-1$ edges and full rank needs at least one root, hence there
are exactly $v-1$ edges and one root. Cycles or excess roots force
column dependence. Conversely, leaf elimination shows that an unweighted
rooted tree has determinant of absolute value one.

Orient every edge of a rooted forest toward its root. Leaf elimination
before row scaling gives the exact product
\begin{equation}\label{eq:forest-product}
 |\det W|=\prod_{\text{nonroot vertices }u}
       (\text{entry at }u\text{ of its parent-edge column}).
\end{equation}
A bad vertex $i\in I$ is never a root; if its parent is $j$, its factor
$b_j/(a_i+b_j)$ is at most $t_i$. A nonroot good vertex has factor
$a_i/(a_i+b_j)<1$. Thus every selected determinant satisfies
\[
 |\det W|\le\prod_{i\in I}t_i=D_*.
\]
Equality holds exactly when every good vertex is a root and every bad
vertex attaches directly to a member of $J_*$. All vertex maxima are
therefore
\begin{equation}\label{eq:vertex-maxima}
 \conv\bigl(\{e_j:j\in J\cup Z\}\cup\{v_{i j_i}:i\in I\}\bigr),
 \qquad j_i\in J_*,
\end{equation}
with independent choices of $j_i$. Ordered by their original bad and
retained indices, all these oriented determinants equal $D_*$. Expanding
an arbitrary simplex's determinant separately in convex combinations of
$K$-vertices proves that its absolute determinant is also at most $D_*$.

\subsection{Every prescribed maximum including nonvertex maxima}
Let $S$ be any maximum and express each of its vertices as a convex
combination of vertices of $K$. Sample the combinations independently.
The oriented determinant of $S$ is the expectation of the sampled
oriented determinant. Equality in
\[
 |\E\det W|\le\E|\det W|\le D_*
\]
forces every positive-probability tuple to be a maximum from
\eqref{eq:vertex-maxima}, with the same orientation.

If no column selected a mandatory retained vertex $e_j$ surely,
independence would give positive probability that the whole tuple omitted
it. Hence each retained vertex occupies a deterministic column. After
removing those columns, each of the disjoint groups
$G_i=\{v_{ij}:j\in J_*\}$ must appear exactly once in every tuple.
The same avoidance argument forces some column to belong surely to each
$G_i$. Distinct groups require distinct columns, accounting for all
remaining columns. Merely relabeling the vertices of this same $S$ gives
\begin{equation}\label{eq:all-maxima}
 \begin{gathered}
 S=\conv\bigl(\{e_j:j\in J\cup Z\}\cup\{w_i:i\in I\}\bigr),\\
 w_i=t_ie_i+s_iz_i,\qquad z_i\in\conv\{e_j:j\in J_*\}.
 \end{gathered}
\end{equation}
Conversely every such simplex has determinant $D_*$ and is maximum.
This includes all continuous tie cases.

\subsection{Exact excess about the same centroid}
Write $z_i=\sum_{j\in J_*}z_{ji}e_j$, with $z_{ji}\ge0$ and
$\sum_jz_{ji}=1$. The $S$-coordinates of an affine point $x$ are
\begin{equation}\label{eq:clip-inverse}
 \begin{aligned}
 \alpha_i(x)&=(1+a_i/B)x_i &&(i\in I),\\
 \alpha_j(x)&=x_j-\sum_{i\in I}(a_i/B)z_{ji}x_i &&(j\in J_*),\\
 \alpha_j(x)&=x_j &&(j\notin I\cup J_*).
 \end{aligned}
\end{equation}
For $j\in J_*$ put $B_{-j}=\max(\{b_h:h\in J,\ h\ne j\}\cup\{0\})$.
The exact largest negative coordinate magnitude is
\begin{equation}\label{eq:clip-beta}
 \beta(S)=\max_{\substack{i\in I\\j\in J_*}}
       \frac{z_{ji}a_iB_{-j}}{B(a_i+B_{-j})}.
\end{equation}
A zero $B_{-j}$ contributes zero. At $v_{ih}$ with $h\ne j$, the possibly
negative coordinate is $-z_{ji}a_ib_h/[B(a_i+b_h)]$. At $v_{ij}$ it is
$s_i(1-z_{ji})\ge0$, and retained vertices have nonnegative coordinates.
Monotonicity of $b/(a_i+b)$ gives \eqref{eq:clip-beta}; convex averaging
covers every $x\in K$. Since $B_{-j}\le B$, it follows that
$\beta(S)\le\max_i a_i/(a_i+B)=\rho_*$.

A centroid dilation by $1+u$ has all simplex coordinates at least $-u/n$.
Consequently the exact identity is
\begin{equation}\label{eq:clip-exact-E}
 E(K,S)=n\beta(S)\le n\rho_*.
\end{equation}
This proves Theorem~\ref{thm:clipping} for the prescribed simplex and its
original centroid.
If $|J_*|\ge2$, then
$\beta(S)=\max_i(s_i\max_{j\in J_*}z_{ji})$. Every vertex maximum has
$E=n\rho_*$, while balanced nonvertex maxima can have smaller excess.
If $J_* =\{j_*\}$, the maximum is unique as a set, and
\[
 \beta(S)=\max_i\frac{a_iB_2}{B(a_i+B_2)},
\]
where $B_2$ is the second-largest positive value, or zero if there is none.
If $|J|=1$, $K$ itself is a simplex and $\beta=0$.
An equal-intercept one-vertex truncation has one bad vertex and, for
$d\ge2$, at least two tied good vertices. Setting $s_i=\rho$ and taking
a vertex maximum gives $\beta=\rho$, proving the stated sharpness.
The graph in this proof concerns clipping vertices; it is not the
facet-normal spanning-two-forest construction for a $d+2$-vertex circuit.

\subsection{An arbitrary body containing an opposite facet}
\begin{theorem}[One-corner bodies]\label{thm:one-corner}
Let $P=\conv(v_0,\ldots,v_d)$ be full dimensional and let $K$ be a convex
body with $\conv(v_1,\ldots,v_d)\subset K\subset P$. Put
$t=\max_K\lambda_0>0$ and $\rho=1-t$. Every prescribed maximum simplex is
\[
 S=\conv(v_1,\ldots,v_d,q),\qquad q\in K,\quad\lambda_0(q)=t,
\]
and $E(K,S)\le(d+1)\rho$. The coefficient is sharp for $d\ge2$ and
$0<\rho<1$. No polyhedral or boundary-regularity hypothesis is imposed.
\end{theorem}
\begin{proof}
If $t=1$, the maximizing point is $v_0$, so $K=P$. Otherwise set
$L=\{x\in P:\lambda_0(x)\le t\}$. By the complete clipping classification,
its maximum volume is $t|P|$, and all maxima consist of the opposite facet
and an apex at height $t$. Compactness gives an apex $q\in K$ attaining
that height, and the entire opposite facet is in $K$. This simplex has
volume $t|P|$, so $K$ and $L$ have the same maximum volume. Every
prescribed maximum of $K$ is therefore a maximum of $L$ and has precisely
the claimed form. Its apex belongs to $K$, without any requirement that
it be an extreme point of $K$.

Write the coordinates of $q$ as $(t,q_1,\ldots,q_d)$, where
$q_j\ge0$ and $\sum_jq_j=\rho$. For $x\in K$, its $S$-coordinates are
\[
 \alpha_0(x)=\lambda_0(x)/t,\qquad
 \alpha_j(x)=\lambda_j(x)-(q_j/t)\lambda_0(x)\ge-q_j\ge-\rho.
\]
The centroid formula proves the stronger estimate
$E(K,S)\le(d+1)\max_jq_j$. Equality in the asserted bound is attained by
$L$ and an apex on one truncated edge, using a second truncated edge to
realize the negative coordinate.
\end{proof}
If the actual enclosure chain supplies a $P$ satisfying the additional
opposite-facet containment hypothesis, the theorem may be used with
$\rho=(H_de(K))^{1/m}$. That extra hypothesis is not asserted for an
arbitrary body or an unrelated enclosing simplex.

\section{The actual defect from signed facet cofactors}\label{sec:cofactor}
Let $K$ be a full-dimensional $d$-polytope with exactly $N=d+2$
irredundant facets. Put $n=d+1$. For facet $i$, let $\nu_i$ be its outward
unit normal, $s_i>0$ its actual area, and $h_i$ its support number. Set
\[
 g_i=s_i\nu_i,\qquad U_i=(g_i,s_ih_i)\in\R^{d+1},\qquad V=|K|>0.
\]
The divergence identities give
\begin{equation}\label{eq:cofactor-sum}
 \sum_i g_i=0,\qquad \sum_i s_ih_i=dV,
 \qquad\sum_iU_i=(0,\ldots,0,dV).
\end{equation}
The horizontal rows of $U$ have rank $d$ because the facet normals span
$\R^d$. The last column-sum identity supplies an independent vertical
direction, so $U$ has rank $n$. Define the signed maximal cofactors
\begin{equation}\label{eq:cofactor-z}
 z_i=(-1)^{i-1}\det(U_1,\ldots,U_{i-1},U_{i+1},\ldots,U_N).
\end{equation}
Then $z\ne0$ and $\sum_i z_iU_i=0$. Translating the origin to the interior
of $K$ leaves these determinants invariant by a row operation and makes
all $h_i>0$. Thus $\sum_i z_is_ih_i=0$ forces both positive and negative
entries in $z$.

\begin{theorem}[Exact cofactor defect]\label{thm:cofactor}
With $L=\sum_i|z_i|$ and $T=\sum_{i<j}|z_i-z_j|$, the original invariant
satisfies
\begin{equation}\label{eq:cofactor-invariant}
 a(K)=\frac LT,\qquad
 \delta(K):=n-\frac1{a(K)}=n-\frac TL,\qquad
 e(K)=\frac{\delta(K)}{n(n-\delta(K))}.
\end{equation}
In particular,
\begin{equation}\label{eq:cofactor-min}
 \delta(K)=\frac2L\sum_{\substack{i<j\\z_iz_j>0}}
                            \min(|z_i|,|z_j|).
\end{equation}
\end{theorem}
\begin{proof}
Let $H=\sum_{i<j}|\det(g_k:k\notin\{i,j\})|$, with the remaining
columns in increasing order. The actual finite-facet identity proved
in Appendix~\ref{prop:actual-moments} gives $a(K)=L/(dVH)$. We show
\begin{equation}\label{eq:cofactor-difference}
 dV\,|\det(g_k:k\notin\{i,j\})|=|z_i-z_j|.
\end{equation}
Fix $i<j$. In the $n$-column matrix omitting column $i$, replace its
column $U_j$ by $\sum_kU_k$. By multilinearity only the original term
and the term replacing $U_j$ by $U_i$ survive. Reordering the latter
columns with the signs in \eqref{eq:cofactor-z} gives $z_i-z_j$ up to an
overall sign. The replacement column is $(0,\ldots,0,dV)$ by
\eqref{eq:cofactor-sum}, and expansion along it gives the other side of
\eqref{eq:cofactor-difference}, again up to sign. Hence $dVH=T$ and
$a=L/T$. Solving $e=a-1/n$ proves \eqref{eq:cofactor-invariant}; all
its denominators are positive.

For real $x,y$, the difference $|x|+|y|-|x-y|$ equals
$2\min(|x|,|y|)$ when $xy>0$, and equals zero otherwise. Each $|z_i|$
occurs in $n$ unordered pairs. Summing this identity gives
\eqref{eq:cofactor-min}.
\end{proof}

\subsection{Signed concentration}
Sort the positive entries as $p_1\ge\cdots\ge p_p>0$ and the absolute
values of the negative entries as $q_1\ge\cdots\ge q_q>0$. Then
\begin{equation}\label{eq:cofactor-sorted}
 \delta(K)=\frac2L\left(\sum_{i=2}^p(i-1)p_i+
                              \sum_{j=2}^q(j-1)q_j\right).
\end{equation}
For $\varepsilon=(L-p_1-q_1)/L$ this gives
\begin{equation}\label{eq:cofactor-concentration}
 2\varepsilon\le\delta(K)\le2\max(p-1,q-1)\varepsilon
                          \le2d\varepsilon.
\end{equation}
Indeed each omitted weight in \eqref{eq:cofactor-sorted} has coefficient
between one and $\max(p-1,q-1)$, and both sign classes are nonempty.
At least $1-\delta(K)/2$ of the normalized absolute cofactor mass is
therefore concentrated on one positive and one negative entry. The
weighted identity, rather than only this comparison, determines the defect.
Common nonzero rescaling and overall sign reversal leave all formulas
unchanged. Zero cofactors, including pyramid degeneracies, are permitted;
no simplicity or genericity condition is required. Arbitrary signed
vectors are not asserted to be realizable facet data.

\subsection{Explicit clipping coordinates}
Suppose $K=P\cap\{f\ge0\}$ with $f=\sum_if_i\lambda_i$, and now assume
that all $n$ old facets and the cut facet remain irredundant. Use the
partition $f_i=-a_i<0$ on $I$, $f_j=b_j>0$ on $J$, allowing zero values.
For any origin there are unique positive scales $r_i,q$ such that
\[
 U_i=r_i(-\nabla\lambda_i,\lambda_i(0)),\qquad
 U_c=q(-\nabla f,f(0)).
\]
Since $\sum_i\lambda_i=1$, the column-sum identity and independence of the
affine barycentric functions imply
\begin{equation}\label{eq:cofactor-scales}
 r_i+qf_i=R,\qquad R=dV>0,
\end{equation}
so each $r_i=R-qf_i$ is positive. The affine identity for $f$ gives the
kernel vector, up to common scale,
\[
 z_c=1,\qquad z_i=-qf_i/r_i.
\]
Write
\[
 \alpha_i=\frac{qa_i}{R+qa_i}\in(0,1),\qquad
 \beta_j=\frac{qb_j}{R-qb_j}>0.
\]
The cut entry and each $\alpha_i$ have the same sign, with
$\min(1,\alpha_i)=\alpha_i$, so \eqref{eq:cofactor-min} becomes
\begin{equation}\label{eq:cofactor-clipping}
 \delta(K)=
 \frac{2\left[\sum_i\alpha_i+\sum_{i<i'}\min(\alpha_i,\alpha_{i'})+
                         \sum_{j<j'}\min(\beta_j,\beta_{j'})\right]}
 {1+\sum_i\alpha_i+\sum_j\beta_j}.
\end{equation}
Zero $f_i$ give zero cofactors. If clipping removes an entire old facet,
that redundant facet must be deleted before applying this calculation;
one may not divide by a vanished area.

For the equal-intercept truncation
\begin{equation}\label{eq:truncation-body}
 K_t=\{x_i\ge0:\ t\le\textstyle\sum_i x_i\le1\},\qquad0<t<1,
\end{equation}
the normalized cofactor vector is
\[
 z_c=1,\qquad z_0=t^d,\qquad z_1=\cdots=z_d=-\beta,
 \qquad\beta=\frac{t^{d-1}(1-t)}{1-t^{d-1}}.
\]
The old opposite facet retains its whole area, each other old facet
retains the fraction $1-t^{d-1}$, and the cut area scale is the fraction
$t^{d-1}$. These facts, or substitution in \eqref{eq:cofactor-scales},
give the displayed vector. Therefore
\begin{equation}\label{eq:truncation-delta}
 \delta(K_t)=\frac{2[t^d+\binom d2\beta]}{1+t^d+d\beta}.
\end{equation}
Substitution into \eqref{eq:cofactor-invariant} yields the exact original
defect and its fixed-dimensional leading term
\begin{equation}\label{eq:truncation-exact}
 \begin{split}
 e(K_t)&=\frac{t^{d-1}[d(d-1)-(d+1)(d-2)t-2t^d]}
 {(d+1)(1-t^d)(d+1+(d-1)t^{d-1})},\\
 e(K_t)&=\frac{d(d-1)}{(d+1)^2}t^{d-1}+O(t^d).
 \end{split}
\end{equation}
For a square followed by $d-2$ pyramids, the four nonzero cofactors are
proportional to $(1,-1,1,-1)$ and the remaining $d-2$ vanish. Thus
$\delta=1$ and $e=1/[d(d+1)]$. The direct facet certificate in the next
section checks the same actual invariant without assuming realizability
of a merely formal cofactor vector.

\section{The asymptotically optimal coefficient for few facets}\label{sec:few}
\subsection{From the actual facet law to a single clipping}
Proposition~\ref{prop:enclosure}, proved in full in the appendices, gives
the precise finite-facet interface needed here: for
$0\le e(K)\le H_d^{-1}$ there is an enclosing simplex $P$, formed by
exactly $d+1$ distinct irredundant $K$-facet halfspaces, such that
\begin{equation}\label{eq:facet-proj}
 \pi_P(u)-\pi_K(u)\le H_de(K)\pi_P(u)\quad(\norm u=1).
\end{equation}
It uses the actual atoms $X_F=\nu_F/h_F$ and probabilities
$p_F=s_Fh_F/(d|K|)$. In particular, the selected halfspaces are actual
facets, rather than unspecified support planes inferred from a compact
limit. The determinant identity is exactly
$(B-A)/B=ne/(1+ne)$ for this law. The assignment bound is
$h\le n^3(n+1)e/(1+ne)$; at the upper endpoint
$h\le1/(3d^2)<1/(2d)$, so its strict enclosure gate remains valid.

We also record the projection-cap argument explicitly.
\begin{lemma}[Near corners from all projections]\label{lem:cap}
If $K\subset P$, with $P$ a full-dimensional simplex, and
$\pi_P(u)-\pi_K(u)\le\delta\pi_P(u)$ for every unit $u$, where
$0\le\delta\le1$, then each corner loss satisfies
\[
 1-\max_{x\in K}\lambda_i^P(x)\le\delta^{1/(d-1)}.
\]
\end{lemma}
\begin{proof}
Let $q_i$ be the corresponding vertex and let
$t_i=1-\max_K\lambda_i^P\in[0,1]$. Choose a unit direction parallel to
its opposite facet. The affine function $\lambda_i^P$ descends under
this projection, and has range $[0,1]$ on the projected simplex.
For $0<s<t_i$, the projected homothetic cap $q_i+s(P-q_i)$ is contained
in $\{\lambda_i^P\ge1-s\}$ and disjoint from the projection of $K$.
Its projected volume is $s^{d-1}\pi_P(u)$. Both sets lie in the
projection of $P$, so $s^{d-1}\le\delta$. Letting $s\uparrow t_i$
proves the claim; $t_i=0$ is immediate. Strict caps avoid any assumption
about the measure of the boundary, and include $\delta=0$.
\end{proof}

\begin{proof}[Proof of the upper bound in Theorem~\ref{thm:few-facet}]
A full-dimensional polytope has at least $d+1$ facets. With exactly $d+1$
it is a simplex, and $E=e=0$. Otherwise it has $d+2$ facets. The $P$
above uses $d+1$ of them, so $K=P\cap H$ for the remaining facet
halfspace. Set $\rho=(H_de(K))^{1/m}$. Lemma~\ref{lem:cap} bounds all
corner losses by $\rho$, and Theorem~\ref{thm:clipping} gives
$E(K,S)\le n\rho$ for this same originally prescribed maximum.
The argument includes $e=0$ and $e=H_d^{-1}$. For larger $e$, the universal
maximum-simplex bound $E\le n$ suffices, since $\rho>1$.

For the uniform scalar estimate, $H_3=8640\le144^2$, and
$H_{d+1}/H_d\le18\le64$ follows from
$(d+1)/d\le2$, $(d+2)/(d+1)\le3/2$, and
$(d+3)/(d+2)\le4/3$ in its squared, cubed, and linear factors.
Induction gives $H_d\le144^{d-1}$. Since $n\le4d/3$,
$nH_d^{1/m}\le192d$.
\end{proof}

\subsection{Sharp exponent and a realized lower coefficient}
The body $K_t$ in \eqref{eq:truncation-body} has exactly $d+2$ facets.
Theorem~\ref{thm:clipping} identifies vertex maxima with
$E(K_t,S_t)=nt$, and \eqref{eq:truncation-exact} gives
$e(K_t)\sim d(d-1)t^m/n^2$. Thus for every fixed $d$ and every
$\eta>1/m$, the ratio $E/e^\eta$ diverges as $t\downarrow0$.
No such stronger exponent can hold uniformly even with a
dimension-dependent coefficient.

For the all-scale lower coefficient take
\[
 C_d=\conv(0,e_1,e_2,e_1+e_2,e_3,\ldots,e_d),\qquad
 S_d=\conv(0,e_1,\ldots,e_d).
\]
This is a square followed by $d-2$ pyramids, with irredundant inequalities
\[
 x_i\ge0\ (1\le i\le d),\qquad
 x_1+\sum_{i=3}^dx_i\le1,\qquad x_2+\sum_{i=3}^dx_i\le1.
\]
Every full-dimensional vertex simplex contains all pyramid vertices and
three of the four square vertices. Every such square triangle has area
$1/2$, so the resulting volume is $1/d!$. Separate determinant expansion
in convex combinations bounds arbitrary inscribed simplices by the same
volume. Thus $S_d$ is maximum. The omitted square vertex has coordinates
$(-1,1,1,0,\ldots,0)$ in $S_d$, so $E(C_d,S_d)=n$.

For completeness, a direct actual-facet certificate gives the defect.
Set $c=1/(d-1)!$. The area-normal/support pairs are
\[
 (-ce_1,0),\quad(-ce_2,0),\quad(-2ce_i,0)\ (3\le i\le d),
\]
\[
 \left(c\left(e_1+\sum_{i=3}^de_i\right),c\right),\qquad
 \left(c\left(e_2+\sum_{i=3}^de_i\right),c\right).
\]
The volume is $2/d!$. For the raw sums of
Appendix~\ref{prop:actual-moments}, direct minors give
\[
 \mathcal H(C_d)=c^d d2^{d-1},\qquad
 \mathcal L(C_d)=c^{d+1}2^d.
\]
Only four lifted subsets contribute, each by $c^{d+1}2^{d-2}$:
omit one top facet or one of the first two coordinate facets.
For horizontal minors the contributions with zero, one, and two top
facets are respectively $1,d,d-1$, multiplied by $c^d2^{d-2}$.
The exact actual identity $a=\mathcal L/(d|C_d|\mathcal H)$ therefore
gives $a(C_d)=1/d$ and $e(C_d)=1/(dn)$. Hence
$G_d^{\rm fac}\ge n[dn]^{1/m}$.

Finally, $H_d$ is a fixed polynomial in $d$, so $H_d^{1/m}\to1$;
also $[dn]^{1/m}\to1$. Divide the two coefficient bounds by $d$ and
squeeze to obtain \eqref{eq:few-squeeze}. No uniformity in a
fixed-dimensional truncation remainder is needed.

\section{Structure forced by all directional projections}
\label{sec:projection-structure}
\label{sec:projection}

Let $P=\conv(p_0,\ldots,p_d)\subset\R^d$ be a full-dimensional
simplex, let $K\subset P$ be a convex body, and assume throughout this
section that
\begin{equation}
 \pi_K(v)\geq (1-\delta)\pi_P(v)
 \quad\text{for every unit direction }v,
 \qquad 0\leq\delta\leq1,\quad d\geq3.
 \label{eq:projection-all-directions}
\end{equation}
Here $\pi_L(v)=\vol_{d-1}(\operatorname{proj}_{v^\perp}L)$ denotes
actual Euclidean projection volume for unit $v$.  Directions constructed below
are normalized when applying this hypothesis.  Write
\[
 n=d+1,\qquad m=d-1,\qquad \rho=\delta^{1/(d-1)},
 \qquad c_P=\frac1n\sum_{i=0}^d p_i.
\]
We first retain an elementary concentric-core argument and its
Brunn--Minkowski refinement, and then prove a stronger finite-spectrum
bound.  These conclusions use the complete hypothesis
\eqref{eq:projection-all-directions}, rather than only the existence of
one point close to each vertex.  They construct a comparison simplex
with centroid $c_P$; this simplex is not asserted to be maximum-volume
and need not coincide with an originally prescribed maximum simplex $S$.

\subsection{Centroid caps and a separating projection}

\begin{lemma}[Elementary centroid cap]
\label{lem:projection-elementary-cap}
Let $D$ be an $m$-dimensional simplex with centroid $g$.  Every closed
halfspace $H$ whose boundary contains $g$ satisfies
\begin{equation}
 \vol_m(D\cap H)\geq \frac{\vol_m(D)}{m+1}.
 \label{eq:projection-elementary-cap}
\end{equation}
Consequently, if $z\in D$, $0\leq\theta\leq1$, and
$y=(1-\theta)z+\theta g$, every closed halfspace $H_y$ through $y$
satisfies
\begin{equation}
 \vol_m(D\cap H_y)\geq
 \frac{\theta^m}{m+1}\vol_m(D).
 \label{eq:projection-homothetic-cap}
\end{equation}
\end{lemma}

\begin{proof}
Use barycentric coordinates $x_0,\ldots,x_m$ on $D$ and write
$H=\{L\geq0\}$ with $L(g)=0$.  Let $\sigma$ cyclically permute these
coordinates.  Each $\sigma^k$ preserves $D$ and its volume, while
\[
 \sum_{k=0}^m L(\sigma^k x)=(m+1)L(g)=0.
\]
At least one summand is nonnegative.  The $m+1$ inverse cyclic images
of $D\cap H$ therefore cover $D$, proving
\eqref{eq:projection-elementary-cap}.  A constant functional causes
no difficulty.  For the second assertion, take the parallel centroid
halfspace and apply the homothety
$x\mapsto(1-\theta)z+\theta x$.  Its image lies in $D\cap H_y$ and
has the volume asserted in \eqref{eq:projection-homothetic-cap}.
\end{proof}

\begin{lemma}[A separator-parallel simplicial projection]
\label{lem:projection-separator}
Let $L$ be a nonconstant affine functional on $P$.  There is a nonzero
direction $v$ parallel to the level hyperplanes of $L$ for which
$D=\operatorname{proj}_{v^\perp}P$ is an $(d-1)$-simplex.  If $g$ is
the centroid of $D$, then
\begin{equation}
 c_P'=\frac d n g+\frac1n r'
 \label{eq:projection-centroid-decomposition}
\end{equation}
for a point $r'\in D$, where primes denote this orthogonal projection.
The functional $L$ descends to a nonconstant affine functional on $D$.
\end{lemma}

\begin{proof}
Choose distinct indices $a,b$ such that $L(p_a)$ and $L(p_b)$ are,
respectively, the minimum and maximum vertex values.  Choose any
$j\notin\{a,b\}$.  A point $r\in[p_a,p_b]$ satisfies
$L(r)=L(p_j)$.  Put $v=p_j-r$.  This vector is nonzero and parallel
to the level hyperplanes of $L$.  Orthogonal projection along $v$
sends $p_j$ to $r'$.  The other $d$ projected vertices are affinely
independent: otherwise $v$ would be parallel to the facet opposite
$p_j$, whereas $p_j-r$ is transverse to that facet.  Their convex hull
is $D$, because $r'$ lies on their projected edge $[p_a',p_b']$.
Summing the projected vertices gives
\eqref{eq:projection-centroid-decomposition}.  The two extreme values
of $L$ are distinct and survive the projection, so the descended
functional is nonconstant.  Repeated extreme values or equality of
$L(p_j)$ with an endpoint value do not affect the construction.
The construction itself requires only $d\geq2$.
\end{proof}

\subsection{Two concentric-core bounds}

\begin{theorem}[Elementary concentric core]
\label{thm:projection-elementary-core}
Define
\[
 \kappa_d=\frac{d+1}{d}\,d^{1/(d-1)}.
\]
If $t=\kappa_d\rho\leq1$, then
\begin{equation}
 (1-t)P+t c_P\subset K.
 \label{eq:projection-elementary-core}
\end{equation}
Moreover,
$\kappa_d\leq4\sqrt3/3<7/3$ for all $d\geq3$.
\end{theorem}

\begin{proof}
Fix $0<t\leq1$ and suppose that
$y=(1-t)x+t c_P\notin K$ for some $x\in P$.  Strict separation gives
a nonconstant affine functional $L$ with
$\max_K L<L(y)$.  Apply Lemma~\ref{lem:projection-separator}, so
that the projection is parallel to the levels of $L$.  By
\eqref{eq:projection-centroid-decomposition},
\begin{align*}
 y'&=(1-t)x'+\frac t n r'+\frac{td}{n}g
     =(1-\theta)z+\theta g,\\
 \theta&=\frac{td}{n},\qquad
 z=\frac{(1-t)x'+(t/n)r'}{1-td/n}\in D.
\end{align*}
The denominator is positive even for $t=1$.  Since $\theta>0$,
$y'$ is interior to $D$ in its affine hull.  The projection of $K$
lies strictly on the lower-$L$ side of $y'$.  The opposite cap has
volume at least
\[
 \frac1d\left(\frac{td}{n}\right)^{d-1}\vol_{d-1}(D)
\]
by Lemma~\ref{lem:projection-elementary-cap}.  In fact, the missing
volume is strictly larger: the strip between $\max_K L$ and $L(y')$
has positive $(d-1)$-volume because $y'$ is interior.  Hence
\begin{equation}
 \delta>\frac1d\left(\frac{td}{n}\right)^{d-1}.
 \label{eq:projection-elementary-contradiction}
\end{equation}
For $t=\kappa_d\rho>0$, the right-hand side is
$\rho^{d-1}=\delta$, a contradiction.  This also proves the
$t=1$ endpoint.  If $\delta=0$, the same contradiction applies for
every $t>0$; letting $t$ decrease to zero and using closedness of $K$
gives $K=P$.

Finally, $d\leq3^{(d-1)/2}$ for $d\geq3$.  There is equality at
$d=3$, and the ratio $(d+1)/d$ is at most $4/3<\sqrt3$ in the
inductive step.  Thus $d^{1/(d-1)}\leq\sqrt3$, and the stated
uniform bound follows from $(d+1)/d\leq4/3$.
\end{proof}

The preceding proof is independent of Brunn--Minkowski.  For completeness
we also retain the stronger centroid-cap argument, including the
one-dimensional derivation used in the refinement.

\begin{lemma}[Centroid cap from Brunn--Minkowski]
\label{lem:projection-sharp-cap}
Let $D$ be an $m$-dimensional convex body with centroid $g$.  Every
closed halfspace $H$ through $g$ satisfies
\begin{equation}
 \frac{\vol_m(D\cap H)}{\vol_m(D)}
 \geq\left(\frac m{m+1}\right)^m.
 \label{eq:projection-sharp-cap}
\end{equation}
\end{lemma}

\begin{proof}
Let $X$ be uniform in $D$ and let $Z=L(X)$, where $L$ is the affine
functional defining $H$, oriented so that $H=\{L\geq0\}$ and
normalized by $\mathbb E Z=0$.  A constant functional is immediate,
so assume $L$ is nonconstant.  Set
\[
 F(u)=\mathbb P(Z\geq u),\qquad R(u)=F(u)^{1/m}.
\]
For $0\leq\alpha\leq1$,
\[
 (1-\alpha)(D\cap\{L\geq u\})+
 \alpha(D\cap\{L\geq v\})
 \subset D\cap\{L\geq(1-\alpha)u+\alpha v\}.
\]
Brunn--Minkowski therefore makes $R$ concave.  Between the minimum and
maximum of $L$ on $D$, this function is continuous and strictly
decreasing.  Its inverse $T$ on $(0,1)$ is also concave and decreasing:
apply concavity of $R$ to $T(s_0),T(s_1)$ and then use its decreasing
monotonicity.  The distribution of $Z$ is continuous, and the
probability-integral transform makes $F(Z)$ uniform on $(0,1)$.
Consequently $U=R(Z)$ has density $m s^{m-1}$ and
\[
 \mathbb E U=\frac m{m+1},\qquad
 \mathbb E T(U)=\mathbb E Z=0.
\]
Jensen's inequality gives $T(m/(m+1))\geq0$.  Since $T$ is
decreasing, $R(0)\geq m/(m+1)$, proving
\eqref{eq:projection-sharp-cap}.  These expectations are bounded
one-dimensional integrals; equivalently, the Jensen step follows by
integrating a supporting line for $T$ at $m/(m+1)$.
\end{proof}

\begin{theorem}[Sharper concentric core]
\label{thm:projection-sharp-core}
The conclusion of Theorem~\ref{thm:projection-elementary-core} holds
with
\begin{equation}
 t=\frac{d+1}{d-1}\rho\leq1.
 \label{eq:projection-sharp-parameter}
\end{equation}
In particular, the coefficient $(d+1)/(d-1)$ is at most $2$.
\end{theorem}

\begin{proof}
Repeat the proof of Theorem~\ref{thm:projection-elementary-core},
replacing the centroid fraction $1/d$ by
$((d-1)/d)^{d-1}$ from Lemma~\ref{lem:projection-sharp-cap}.
The same homothety and strictly positive strip show that a missing
core point would force
\[
 \delta>
 \left(\frac{t(d-1)}{d+1}\right)^{d-1}.
\]
The choice \eqref{eq:projection-sharp-parameter} contradicts this
inequality.  Strict separation again includes $t=1$, and the closed
limit used above includes $\delta=0$.
\end{proof}

\subsection{An elementary finite-spectrum refinement}

\begin{theorem}[Finite-spectrum concentric core]
\label{thm:projection-finite-spectrum}
Put
\[
 M_0=1,\qquad M_k=\frac{k^k}{(k+1)^{2k+1}}\quad(k\geq1),
 \qquad
 \tau_d(\delta)=n\max_{1\leq\ell\leq m}
                   (M_{\ell-1}\delta)^{1/\ell}.
\]
If $0\leq\tau_d(\delta)\leq1$, then
\begin{equation}
 (1-\tau_d(\delta))P+\tau_d(\delta)c_P\subset K.
 \label{eq:projection-finite-spectrum-core}
\end{equation}
This conclusion uses neither Brunn--Minkowski nor a centroid-cap
inequality.  One may use the smallest parameter among the applicable
core bounds.
\end{theorem}

\begin{proof}
We first obtain an explicit missing-cap estimate.  Let $b\geq0$ be
affine on $P$, normalized so that its minimum vertex value is zero
and the sum of all its vertex values is one.  List its other $d$
vertex values as
\[
 B\geq C\geq b_1\geq\cdots\geq b_{d-2}\geq0.
\]
Let $0<s\leq1/n$ and assume $K$ lies strictly in $\{b>s\}$.
Project along the line joining the vertex of value $C$ to the point
of value $C$ on the edge with endpoint values $0$ and $B$.
As in Lemma~\ref{lem:projection-separator}, this direction is
parallel to the levels of $b$ and gives an $m$-simplex $D$ whose
vertex values are
\[
 0,B,b_1,\ldots,b_{d-2}.
\]
The value-$C$ vertex has become redundant.  Since $B\geq1/d>s$,
the cap $\{b\leq s\}$ contains the simplex with apex at the
zero-valued vertex and edge-intercept fractions
$\min(1,s/a)$ for each other endpoint value $a$, with fraction one
when $a=0$.  Its relative volume is exactly
\begin{equation}
 R_s=\frac{s^m}{B\prod_{i=1}^{d-2}\max(b_i,s)}.
 \label{eq:projection-intercept-volume}
\end{equation}
All of it is missing from the projection of $K$.  The actual missing
volume is strictly larger, since $K$ lies strictly above level $s$
and $D$ has positive-volume strips immediately above that level.

Let $k$ be the number of $b_i$ exceeding $s$, so $0\leq k\leq m-1$.
For $k=0$ the denominator in \eqref{eq:projection-intercept-volume}
is at most $s^{m-1}=M_0s^{m-1}$.  For $k\geq1$, let $U$ be the sum
of these $k$ active values.  As $C$ is at least each active value,
$C\geq U/k$, and normalization gives
\[
 1\geq B+C+U\geq B+(1+1/k)U.
\]
The arithmetic--geometric mean inequality yields
\[
 \prod_{\mathrm{active}}b_i
 \leq(U/k)^k\leq\left(\frac{1-B}{k+1}\right)^k.
\]
The maximum of $B((1-B)/(k+1))^k$ on $[0,1]$ occurs at
$B=1/(k+1)$ and equals $M_k$.  Hence
\[
 B\prod_i\max(b_i,s)
 \leq\max_{0\leq k\leq m-1}M_k s^{m-1-k}.
\]
The relative projection deficit is therefore strictly larger than
\begin{equation}
 \min_{1\leq\ell\leq m}\frac{s^\ell}{M_{\ell-1}}.
 \label{eq:projection-optimized-cap}
\end{equation}
This optimization is an upper bound on the denominator; it does not
assert that every formal optimizer is compatible with every active
coordinate regime.

Now suppose $y=(1-t)x+t c_P\notin K$, with $x\in P$ and
$0<t\leq1$.  Orient a strict separating affine functional so that
$K$ has larger values than $y$, subtract its minimum vertex value,
and divide by the sum of its vertex values.  This gives the preceding
normalized $b$, with
\[
 b(y)=(1-t)b(x)+t/n\geq t/n.
\]
Thus $K$ lies strictly in $\{b>t/n\}$.  Set $s=t/n$ in
\eqref{eq:projection-optimized-cap}.  If $t=\tau_d(\delta)>0$,
then every term $s^\ell/M_{\ell-1}$ is at least $\delta$, a
contradiction to \eqref{eq:projection-all-directions}.  The argument
includes $t=1$.  For $\delta=0$, apply the positive-$t$ argument
and pass to the closed limit as before.
\end{proof}

\begin{proposition}[Uniform top-spectrum branch]
\label{prop:projection-top-spectrum}
For every $d\geq3$ and $0\leq\rho\leq1/3$, the term $\ell=m$ is
largest in the definition of $\tau_d$.  Consequently
\begin{align}
 \tau_d(\delta)&=\gamma_d\rho,\nonumber\\
 \gamma_d&=n\left(
       \frac{(d-2)^{d-2}}{(d-1)^{2d-3}}\right)^{1/(d-1)}
   =\frac{m+2}{m}\left(1-\frac1m\right)^{(m-1)/m}
   \longrightarrow1.
 \label{eq:projection-gamma}
\end{align}
On this branch $\tau_d(\delta)\leq2/3<1$, so the positive-core
condition is automatic.  In each fixed dimension the exact threshold
\begin{equation}
 \delta\leq
 \min_{1\leq\ell<m}
 \left(\frac{M_{m-1}^{\ell}}{M_{\ell-1}^{m}}
 \right)^{1/(m-\ell)}
 \label{eq:projection-exact-branch}
\end{equation}
is also sufficient and may enlarge this branch.  In particular,
\[
 \tau_3(\delta)=\max(4\delta,\sqrt{2\delta}),
\]
which equals $\sqrt2\,\rho$ for $\delta\leq1/8$.
\end{proposition}

\begin{proof}
Set $f_\ell=\ell M_{\ell-1}^{1/\ell}$.  Then $f_1=1$ and
\[
 f_\ell=(1-1/\ell)^{1-1/\ell}\quad(\ell\geq2).
\]
The latter sequence increases, since $x\log x$ increases on
$[1/2,1)$.  In particular, $1/\sqrt2\leq f_m<1$.
For $\rho>0$ and $2\leq\ell\leq m$, the ratio of the $\ell$th
term to the $m$th term is
\[
 \frac m\ell\frac{f_\ell}{f_m}\rho^{m/\ell-1}
 \leq\frac m\ell\rho^{m/\ell-1}\leq1
 \qquad(\rho\leq1/3).
\]
For the last inequality put $x=m/\ell\geq1$ and use
$\log x\leq x-1\leq(x-1)\log3$.  For $\ell=1$ the ratio is
\[
 \frac m{f_m}\rho^{m-1}
 \leq\sqrt2\,m3^{1-m}\leq\frac{2\sqrt2}{3}<1.
\]
Indeed $m3^{1-m}$ decreases for integer $m\geq2$, because its
successive ratio is $(m+1)/(3m)<1$.  The case $\rho=0$ is immediate
without division.  Formula \eqref{eq:projection-gamma} follows by
substitution, and
\[
 \tau_d(\delta)=\frac{m+2}{m}f_m\rho\leq\frac23<1.
\]
For fixed $m$ and $\ell<m$, the ratio
$(\delta M_{\ell-1})^{1/\ell}/
(\delta M_{m-1})^{1/m}$ tends to zero as $\delta$ decreases to
zero.  Rearranging the condition that this ratio be at most one for
every $\ell<m$ gives \eqref{eq:projection-exact-branch}.
The $d=3$ formula follows from $M_0=1$ and $M_1=1/8$.
\end{proof}

\subsection{Depth at faces of every positive dimension}

Let $\lambda_0,\ldots,\lambda_d$ be the barycentric coordinates of
$P$.  For a nonempty index set $I$ define
$f_I=\sum_{i\in I}\lambda_i$ and $s_I=\min_K f_I$.

\begin{theorem}[Face-depth hierarchy]
\label{thm:projection-face-depth}
If $|I|=r$ with $1\leq r\leq d-1$, then
\begin{equation}
 F_{r,d-r}(s_I)\leq\delta,\qquad
 F_{r,d-r}(s)=\sum_{j=r}^{d-1}
     \binom{d-1}{j}s^j(1-s)^{d-1-j}.
 \label{eq:projection-face-polynomial}
\end{equation}
In particular,
\begin{equation}
 s_I\leq\delta^{1/r}=\rho^{(d-1)/r}.
 \label{eq:projection-face-power}
\end{equation}
For each facet coordinate,
\begin{equation}
 \min_K\lambda_i\leq1-(1-\delta)^{1/(d-1)}
 \leq\frac{2\delta}{d-1}\quad\text{if }\delta\leq\tfrac12.
 \label{eq:projection-facet-depth}
\end{equation}
\end{theorem}

\begin{proof}
The complement of $I$ contains two indices $a,b$.  Project along
the edge direction $p_a-p_b$.  The image $D$ is an $(d-1)$-simplex,
obtained by identifying exactly the vertices $a,b$.  The affine
function $f_I$ descends and becomes the sum of exactly $r$ of the
$d$ barycentric coordinates of $D$.  Since $K\subset\{f_I\geq s_I\}$,
its projection omits $\{f_I<s_I\}$.

A direct section calculation in a simplex with $d$ vertices gives
the density of the sum of $r$ barycentric coordinates as
\[
 \frac{(d-1)!}{(r-1)!(d-r-1)!}
 s^{r-1}(1-s)^{d-r-1},\qquad 0<s<1.
\]
Its integral from zero to $s$ is $F_{r,d-r}(s)$: differentiating
the polynomial in \eqref{eq:projection-face-polynomial} makes
adjacent binomial terms telescope to the displayed density, and the
polynomial vanishes at zero.  No independence of barycentric
coordinates is asserted.  The omitted relative volume is therefore
$F_{r,d-r}(s_I)$, proving the first inequality.

As a polynomial identity, $F_{r,d-r}(s)$ also equals the probability
of at least $r$ successes among $d-1$ independent Bernoulli trials
with parameter $s$.  The event that the first $r$ trials succeed
has probability $s^r$ and is contained in that event.  Thus
$F_{r,d-r}(s)\geq s^r$, proving \eqref{eq:projection-face-power}.
For $r=1$, $F_{1,d-1}(s)=1-(1-s)^{d-1}$, which gives the exact
first bound in \eqref{eq:projection-facet-depth}.  For
$m=d-1$ and $\delta\leq1/2$,
\[
 1-(1-\delta)^{1/m}
 =\int_0^\delta\frac1m(1-u)^{1/m-1}\,du
 \leq\frac{\delta}{m(1-\delta)}\leq\frac{2\delta}{m}.
\]
The zero and unit depth endpoints follow directly or by continuity.
\end{proof}

The theorem concerns faces of positive dimension: the face
$\{\lambda_i=0:i\in I\}$ has codimension $r\leq d-1$.
The vertex case is separate, with projection exponent $d-1$, rather
than $d$.  When $\rho$ is of order $1/d$, the facet estimate imposes
exponentially small barycentric distance to each facet.

\subsection{Necessary tests for finite hulls and the remaining target}

Let $K_0$ be the convex hull of finitely many proposed near-corner
points and candidate-simplex vertices.  Since $f_I$ is affine,
$\min_{K_0}f_I$ is the minimum over the generators.  Thus a violation
of \eqref{eq:projection-face-polynomial} rules out
\eqref{eq:projection-all-directions}, regardless of any replacement
minor constraints and without deciding whether the candidate simplex
is maximum-volume.

\begin{proposition}[A facet obstruction for finite hulls]
\label{prop:projection-finite-hull-obstruction}
If every generator of $K_0$ has $\lambda_i\geq\varepsilon$ for
some fixed $i$ and $0\leq\varepsilon\leq1$, then any deficit
$\delta$ satisfying \eqref{eq:projection-all-directions} obeys
\begin{equation}
 \delta\geq1-(1-\varepsilon)^{d-1}
 \geq\frac{(d-1)\varepsilon}{1+(d-1)\varepsilon}.
 \label{eq:projection-facet-obstruction}
\end{equation}
\end{proposition}

\begin{proof}
The first inequality is the facet case of
Theorem~\ref{thm:projection-face-depth}.  For $\varepsilon<1$,
$(1-\varepsilon)^{-(d-1)}\geq1+(d-1)\varepsilon$, which gives
the second inequality after taking reciprocals.  The endpoint
$\varepsilon=1$ follows directly.
\end{proof}

For example, let the comparison simplex have barycentric matrix
$Q_a=(1-a)\Id_n+(a/n)\ones\ones^{\mathsf T}$, and suppose every
additional generator has every coordinate at least $b/n$.
Then \eqref{eq:projection-facet-obstruction} applies to every facet
with $\varepsilon=\min(a,b)/n$.  If $a$ and $b$ are comparable to
$\rho$, the actual deficit is bounded below by a constant times
$\rho$, rather than $\rho^{d-1}$.  Such thin positive-coordinate
hulls cannot model the difficult $\rho\asymp1/d$ regime.

The core theorem gives another finite test: whenever
$\kappa_d\rho\leq1$, every point
\[
 y_i=(1-\kappa_d\rho)p_i+\kappa_d\rho c_P
\]
must belong to $K_0$.  Membership is a finite linear feasibility
problem.  If membership fails, its separating functional determines,
by Lemma~\ref{lem:projection-separator}, an actual projection
direction violating the proposed deficit.  This is only a necessary
test; passing it does not establish all-direction projection control.
The stronger core parameters give the corresponding stronger tests.

In particular, the full projection hypothesis permits the special
central comparison matrix
\[
 Q_t=(1-t)\Id_n+\frac t n\ones\ones^{\mathsf T}.
\]
For a prescribed maximum simplex with barycentric matrix $W$, maximum
volume still gives all replacement-minor inequalities for
$W^{-1}Q_t$.  A dimension-free lower bound on its entries, or a proof
that the additional face-depth inequalities force such a bound,
remains open here.  These conclusions alone do not improve the
general estimate $E(K,S)\leq16n^2\rho$ to $Cn\rho$ and do not
produce a genuine quadratic lower-bound body.

\begin{corollary}[Concentric comparison without prescribed retention]
\label{cor:projection-concentric-comparison}
Let $t<1$ be any core parameter proved above and put
$T=(1-t)P+t c_P$.  Then $T$ is a full-dimensional simplex with
centroid $c_P$, and
\begin{equation}
 T\subset K\subset P
 =c_P+\frac1{1-t}(T-c_P),\qquad
 E(K,T)\leq\frac t{1-t}.
 \label{eq:projection-comparison-excess}
\end{equation}
For the sharper parameter $t=((d+1)/(d-1))\rho\leq1/2$, this gives
$E(K,T)\leq4\rho$.  On the uniform finite-spectrum branch
$\rho\leq1/3$, it gives
\begin{equation}
 E(K,T)\leq\frac{\gamma_d\rho}{1-\gamma_d\rho},
 \qquad \gamma_d\longrightarrow1.
 \label{eq:projection-finite-excess}
\end{equation}
\end{corollary}

\begin{proof}
The core inclusion and $K\subset P$ give the sandwich, and the
displayed homothety is exact.  The defining centered excess therefore
gives \eqref{eq:projection-comparison-excess}.  If $t\leq1/2$,
then $t/(1-t)\leq2t\leq4\rho$.  The finite-spectrum assertion
follows by substituting $t=\gamma_d\rho$ and using
Proposition~\ref{prop:projection-top-spectrum}.
\end{proof}

\begin{remark}[Actual-defect substitution and scope]
\label{rem:projection-actual-defect-scope}
On the small-defect branch where the previously established actual
enclosing-simplex/projection chain supplies
$\delta=H_de(K)$ with $H_d=3d^2(d+1)^3(d+2)$, one may substitute
\[
 t=\frac{d+1}{d-1}H_d^{1/(d-1)}e(K)^{1/(d-1)}
\]
in \eqref{eq:projection-comparison-excess} whenever $t<1$.
Likewise the finite-spectrum parameter may be applied to the same
actual projection packet, within its stated gate.  Neither
substitution changes the defect or its power law.  The simplex $T$
retains the enclosing centroid $c_P$; it is not identified with an
arbitrarily prescribed maximum simplex $S$, and its centroid is not
substituted for the original centroid of $S$.  Thus the
dimension-free comparison approximation and prescribed-maximum
retention remain distinct assertions.
\end{remark}

\section{Inverse minor rigidity in structured classes}
\label{sec:inverse}

Let $W$ be an invertible nonnegative column-stochastic $n\times n$
matrix, and write
\[
 B=W^{-1},\qquad
 M_k(B)=\max_{|I|=|J|=k}|\det B_{I,J}|,\qquad
 a(B)=\max\{0,-\min_{i,j}B_{ij}\}.
\]
The general problem asks for a dimension-free bound
$a(B)\leq C(c-1)$ when $M_k(B)\leq c^k$ for every $1\leq k\leq n$,
with $c\geq1$.  We prove such bounds for support-two matrices,
low-order blocks, and symmetric positive-definite matrices.  The
first class yields an all-scale geometric estimate for every
prescribed maximum simplex of an edge-supported body.  The
unrestricted inverse-minor problem is not resolved by these results.

\subsection{Weighted cycles}

\begin{lemma}[Weighted cyclic blocks]
\label{lem:inverse-cycles}
Suppose a column-stochastic block of order $q\geq2$ has the form
\[
 W=(\Id_q+P R_d)D,\qquad
 D=\diag(d_i),\quad R_d=\diag(r_i),\quad
 d_i=\frac1{1+r_i},\quad r_i>0,
\]
where $P$ is a cyclic permutation matrix.  If
$M_1(B)\leq c$ and $M_2(B)\leq c^2$, then
\begin{align}
 a(B)&\leq c^2-1 &&(1\leq c<\sqrt2),\label{eq:inverse-cycle-local}\\
 a(B)&\leq(2+\sqrt2)(c-1)&&(c\geq1).
 \label{eq:inverse-cycle-global}
\end{align}
\end{lemma}

\begin{proof}
Let $R=\prod_i r_i$.  Interchanging the two cyclic matchings by a
row permutation replaces this product by $R^{-1}$.  We may choose
the orientation with $R\leq1$.  In even order, $R=1$ is excluded by
invertibility.  Thus
\[
 L=1-(-1)^qR>0.
\]
The finite geometric identity gives
\begin{equation}
 B=\frac1L D^{-1}
 \bigl[\Id_q-PR_d+(PR_d)^2-\cdots+
                   (-1)^{q-1}(PR_d)^{q-1}\bigr].
 \label{eq:inverse-cycle-expansion}
\end{equation}
In particular,
\begin{equation}
 B_{ii}=\frac{1+r_i}{L},\qquad
 \det B_{\{i,j\},\{i,j\}}=
 \frac{(1+r_i)(1+r_j)}{L}\quad(i\ne j).
 \label{eq:inverse-cycle-two-minor}
\end{equation}
For the second formula, Jacobi's complementary-minor identity
applies.  Directly, when $q>2$, deleting vertices $i,j$ from the
directed cycle leaves an acyclic principal submatrix, triangular
after topological ordering, with determinant the product of its
$d_\ell$.  For $q=2$, the formula is simply
$\det B=1/(d_1d_2L)$.

First suppose every $r_i\leq1$.  If $q$ is odd, then $R\leq r_i$,
so $B_{ii}\geq1$; if $q$ is even, $L\leq1$ gives the same
conclusion.  In row $i$, each absolute inverse entry in
\eqref{eq:inverse-cycle-expansion} is the common factor
$(1+r_i)/L$ times a product of consecutive ratios.  Negative entries
correspond to odd-length paths.  Every such product is at most the
last one-edge factor $r_j$, where $j$ is the predecessor of $i$.
Thus the largest negative magnitude in that row occurs on the
incoming edge and equals
\[
 -B_{ij}=\frac{(1+r_i)r_j}{L}.
\]
The principal two-minor on $i,j$ is
\[
 \frac{(1+r_i)(1+r_j)}L
 =B_{ii}+(-B_{ij})\geq1+(-B_{ij}).
\]
It follows that $a(B)\leq M_2(B)-1\leq c^2-1$.

Next suppose some ratio is at least one, and let $x\geq y$ be the
two largest ratios, with $x\geq1$.  Then $R\leq xy$: if $y\leq1$
all the remaining ratios are at most one, while if $y\geq1$ then
$xy\geq1\geq R$.  In odd order,
\begin{align*}
 (1+x)(1+y)-2(1+R)
 &\geq(1+x)(1+R/x)-2(1+R)\\
 &=(x-1)(1-R/x)\geq0.
\end{align*}
The corresponding two-minor in
\eqref{eq:inverse-cycle-two-minor} is therefore at least two.  In
even order $L\leq1$ and $(1+x)(1+y)\geq2$ give the same result.
Hence this case requires $c\geq\sqrt2$.

Consequently $c<\sqrt2$ forces the first case and proves
\eqref{eq:inverse-cycle-local}.  For $c\geq\sqrt2$, the entry bound
$a(B)\leq c$ yields
\[
 \frac{a(B)}{c-1}\leq\frac c{c-1}\leq2+\sqrt2.
\]
For $1<c<\sqrt2$, the first case gives
$a(B)/(c-1)\leq c+1<1+\sqrt2$.  At $c=1$ its conclusion is
$a(B)=0$.  This proves \eqref{eq:inverse-cycle-global}.
\end{proof}

Independent row and column permutations of $W$ merely permute the
rows and columns of $B$ independently and change signs of its minors.
All these inequalities are therefore invariant under such
permutations.  Zero edge weights may be included by continuity, but
the support decomposition below handles them without treating a
zero-weight edge as part of a genuine cycle.

\subsection{Columns supported on at most two rows}

\begin{theorem}[Support-two inverse-minor rigidity]
\label{thm:inverse-support-two}
Suppose every column of $W$ has at most two nonzero entries.  If
$M_1(B)\leq c$ and $M_2(B)\leq c^2$, then
\begin{align}
 a(B)&\leq c^2-1 &&(1\leq c<\sqrt2),
 \label{eq:inverse-support-local}\\
 a(B)&\leq(2+\sqrt2)(c-1)&&(c\geq1).
 \label{eq:inverse-support-global}
\end{align}
For $n=1$ only the order-one bound is needed.  The constants do not
depend on the matrix order, cycle lengths, or sizes of attached trees.
\end{theorem}

\begin{proof}
Associate a multigraph to $W$.  Its vertices are the row indices;
a column with two positive entries is an edge between its two
supporting rows, and a singleton column is a loop at its supporting
row.  Parallel edges are allowed.  Connected components partition
the rows and columns into matrix blocks.  Each block must be square:
if a component has unequal row and column counts, summing the
component rank bounds, each at most the smaller count, makes the
whole square matrix singular.  Every connected component therefore
has equally many edges and vertices and is unicyclic, including
loops and doubled edges as cycles of lengths one and two.

Remove the trees attached to the unique cycle.  Pair each non-core
vertex $v$ with the edge joining it to its parent $p(v)$ towards
the core.  Order core rows and columns first and tree vertices in
parent-before-child order.  Then
\[
 W=\begin{pmatrix}W_C&F\\0&T\end{pmatrix},
\]
where $T$ is triangular.  The column paired with $v$ has weight
$d_v>0$ at $v$ and $1-d_v$ at its parent.  Its full inverse diagonal
entry is $1/d_v$, so
\begin{equation}
 r_v:=\frac{1-d_v}{d_v}=\frac1{d_v}-1\leq c-1.
 \label{eq:inverse-tree-ratio}
\end{equation}
The core inverse $B_C=W_C^{-1}$ is a submatrix of $B$ and inherits
the order-one and order-two bounds.  A loop core has $B_C=[1]$;
the other cores satisfy Lemma~\ref{lem:inverse-cycles}.

For completeness, let $b_v=B e_v$ denote an inverse column.  The
identity $BW e_v=e_v$ gives the exact recursion
\begin{equation}
 b_v=d_v^{-1}e_v-r_v b_{p(v)}.
 \label{eq:inverse-tree-recursion}
\end{equation}
For a core index $p$, $b_p$ has core part $B_Ce_p$ and zero tree
part.  Repeatedly applying \eqref{eq:inverse-tree-recursion} shows
that every off-diagonal entry in a tree column is a product of at
least one tree ratio, times a core inverse entry or an ancestor
diagonal entry $1/d_u$.  Every such starting factor has absolute
value at most $c$.  Equivalently, the tree block has the form
$(\Id+N)D$ with $N_{p(v),v}=r_v$, and the finite nilpotent expansion
of $(\Id+N)^{-1}$ gives these unique path products.

If $c<\sqrt2$, then $0\leq r_v\leq c-1<1$.  Every tree-column
off-diagonal magnitude is therefore at most
$c(c-1)\leq c^2-1$.  Tree diagonal entries are positive, and core
negative entries satisfy the same $c^2-1$ bound by the cyclic
lemma.  This proves \eqref{eq:inverse-support-local}.  For
$c\geq\sqrt2$, the bound $a(B)\leq c$ gives
\eqref{eq:inverse-support-global} as in that lemma; for smaller
$c$ the local bound already implies it.  Taking direct sums and
undoing the permutations completes the proof.
\end{proof}

\subsection{Retention of a prescribed maximum simplex}

The next reduction records why the matrix estimates control the
entire body, including points beyond the selected near-vertex
comparison simplex.

\begin{lemma}[From near-vertex points to inverse minors]
\label{lem:inverse-geometric-reduction}
Let $P$ have $n$ vertices, let $K\subset P$ be a convex body, and
let $S$ be a prescribed maximum-volume simplex in $K$.  Let $W$ be
the matrix of its vertex coordinates in the barycentric system of
$P$.  Suppose $K$ contains $n$ points whose nonnegative
column-stochastic barycentric matrix $Q$ satisfies
$Q_{ii}\geq1-\rho$.  For $0\leq\rho<1/2$, put
\[
 B=W^{-1},\qquad A=BQ,\qquad c=(1-2\rho)^{-1}.
\]
Then every square minor of $A$ has absolute value at most one,
and $M_k(B)\leq c^k$ for every $k$.  Moreover,
\begin{equation}
 E(K,S)\leq n\min\{1,a(B)\}.
 \label{eq:inverse-all-body-excess}
\end{equation}
The latter conclusion applies to all of $K$ and uses the original
centroid of the prescribed $S$.
\end{lemma}

\begin{proof}
Replacing any selected set of vertices of $S$ by selected columns
of $Q$ gives a simplex in $K$.  The determinant ratio is the
corresponding square minor of $A$, up to sign.  Maximality of $S$
therefore gives the all-minor bound.  Since
$\|Q-\Id_n\|_1\leq2\rho$, the Neumann series gives
$\|Q^{-1}\|_1\leq c$.  From $B=AQ^{-1}$, Cauchy--Binet bounds
each $k$-minor of $B$ by the sum of absolute $k$-minors of
$Q^{-1}$ with a fixed column set.  Expanding their determinants
over injections bounds that sum by the product of the corresponding
$k$ absolute column sums, at most $c^k$.

For any $x\in K$, its $P$-barycentric vector $z$ is nonnegative and
sums to one.  Its $S$-barycentric vector is $Bz$, whose every entry
is at least $\min B$.  Single-vertex replacement maximality also
gives $|\alpha_i^S(x)|\leq1$.  The original-centroid dilation
criterion \eqref{eq:dilationcoords} is
$\alpha_i^S(x)\geq-E/n$.  Combining these two lower bounds for
every $x\in K$ proves \eqref{eq:inverse-all-body-excess}.
\end{proof}

\begin{theorem}[Edge-supported prescribed simplex]
\label{thm:inverse-edge-simplex}
Under the hypotheses of Lemma~\ref{lem:inverse-geometric-reduction},
suppose every vertex of the originally prescribed maximum simplex
$S$ lies on an edge of $P$, allowing vertices of $P$ themselves.
For every $\rho\geq0$ for which the stated near-vertex hypothesis
holds,
\begin{equation}
 E(K,S)\leq(4+2\sqrt2)n\rho.
 \label{eq:inverse-edge-geometric}
\end{equation}
No bound on the number of vertices of $K$ is required.
\end{theorem}

\begin{proof}
The support assumption says every column of $W$ has at most two
nonzero entries.  Put
\[
 \rho_0=\frac{1-1/\sqrt2}{2}.
\]
For $0\leq\rho<\rho_0$, the reduction has
$c=(1-2\rho)^{-1}<\sqrt2$.  By
Theorem~\ref{thm:inverse-support-two} and
\eqref{eq:inverse-all-body-excess},
\[
 \frac{E(K,S)}n\leq c^2-1
 =\frac{4\rho(1-\rho)}{(1-2\rho)^2}
 \leq(4+2\sqrt2)\rho.
\]
Indeed $4(1-\rho)/(1-2\rho)^2$ increases on this interval and
equals $4+2\sqrt2$ at $\rho_0$.  For $\rho\geq\rho_0$, universal
single-replacement maximality gives $E(K,S)\leq n$, and
$1/\rho_0=4+2\sqrt2$.  This proves the entire large-$\rho$ branch,
including $\rho\geq1/2$ without using the Neumann reduction.
\end{proof}

\begin{theorem}[Every prescribed maximizer of an edge-supported body]
\label{thm:inverse-edge-body}
Let $K\subset P$ be a full-dimensional polytope whose every vertex
lies on an edge of $P$.  Suppose $K$ contains the near-vertex points
in Lemma~\ref{lem:inverse-geometric-reduction}.  Then
\eqref{eq:inverse-edge-geometric} holds for every prescribed
maximum-volume simplex $S$ of $K$, even if its own vertices are
not edge-supported.
\end{theorem}

\begin{proof}
Let $v_1,\ldots,v_N$ be the vertices of $K$.  Write every vertex
$s_j$ of the prescribed $S$ as a convex combination
$s_j=\sum_\ell\lambda_{\ell j}v_\ell$.  Independently sample
$T_j=v_\ell$ with probabilities $\lambda_{\ell j}$.  Use homogeneous
affine coordinates, in which the determinant of the $n$ columns
is proportional to oriented simplex volume, and order $S$ so that
$\det S=D>0$.  Determinant multilinearity gives
\[
 \mathbb E[\det T]=\det S=D.
\]
Maximality gives $|\det T|\leq D$ for every sampled tuple.  Hence
every positive-probability tuple has $\det T=D$: it is
nondegenerate, has the same orientation, and is a maximum vertex
simplex.

Fix $x\in K$ and an index $j$.  Cramer's rule and multilinearity
in all unreplaced columns give
\begin{align}
 \alpha_j^S(x)
 &=\frac{\det(S\text{ with column }j\text{ replaced by }x)}D
 \nonumber\\
 &=\frac{\mathbb E[\det(T\text{ with column }j
                              \text{ replaced by }x)]}D
  =\mathbb E[\alpha_j^T(x)].
 \label{eq:inverse-maximum-averaging}
\end{align}
Each sampled $T$ is edge-supported.  The coordinate bound in the
proof of Theorem~\ref{thm:inverse-edge-simplex} therefore gives
$\alpha_j^T(x)\geq-(4+2\sqrt2)\rho$ for all $x,j$.
Averaging proves the identical coordinate bound for the original
$S$.  The dilation criterion then yields
\eqref{eq:inverse-edge-geometric} with $S$ and its own centroid.
No sampled centroid is substituted.
\end{proof}

\begin{corollary}[Compact bodies supported on the one-skeleton]
\label{cor:inverse-extreme-skeleton}
Theorem~\ref{thm:inverse-edge-body} also applies when $K$ is a compact
convex body and every extreme point of $K$ lies on the finite
one-skeleton of $P$.
\end{corollary}

\begin{proof}
For every edge $F$ of $P$, $K\cap F$ is empty, a point, or a closed
interval.  A point in the relative interior of a nontrivial such
interval is not extreme in $K$.  Thus each edge contains at most
two extreme points of $K$, its interval endpoints.  There are only
finitely many edges, so $K$ has finitely many extreme points and,
by the finite-dimensional compact convex hull theorem, is their
convex hull.  It is therefore a polytope to which
Theorem~\ref{thm:inverse-edge-body} applies.  In particular, the same
conclusion holds when $K$ is the convex hull of finitely many closed
segments on the edges of $P$.
\end{proof}

\subsection{A conditional corollary for the actual defect}

\begin{corollary}[Actual-defect retention with compatible edge support]
\label{cor:inverse-actual-defect}
Retain the original actual invariant
\begin{equation}
 e(K)=\left(\frac d{d+1}\right)^d
       \frac{R(\pyr K)}{R(K)}-1-\frac1{d+1},
 \qquad
 R(L)=\frac{\vol(\Pi L)}{\vol(L)^{\dim L-1}},
 \label{eq:inverse-original-defect}
\end{equation}
and put $H_d=3d^2(d+1)^3(d+2)$, $n=d+1$, $d\geq3$.
On the branch $0\leq e(K)\leq H_d^{-1}$, assume additionally that
at least one enclosing simplex $P$ admitted by the actual-defect,
projection, and cap chain has every extreme point of $K$ on its
one-skeleton.  Then for every originally prescribed maximum simplex
$S$,
\begin{equation}
 E(K,S)\leq(4+2\sqrt2)(d+1)H_d^{1/(d-1)}
                         e(K)^{1/(d-1)}.
 \label{eq:inverse-actual-edge}
\end{equation}
In particular,
\begin{equation}
 E(K,S)\leq(768+384\sqrt2)d\,e(K)^{1/(d-1)}.
 \label{eq:inverse-actual-edge-coarse}
\end{equation}
For $e(K)>H_d^{-1}$, both bounds hold without the compatibility or
edge-support hypothesis.
\end{corollary}

\begin{proof}
On the stated small-defect branch, the retained actual chain
supplies the compatible enclosing $P$, relative projection deficit
at most $\delta=H_de(K)$, and near-vertex radius
$\rho=\delta^{1/(d-1)}$.  The extreme-point assumption makes $K$
a polytope by Corollary~\ref{cor:inverse-extreme-skeleton}.
Substitution into Theorem~\ref{thm:inverse-edge-body} gives
\eqref{eq:inverse-actual-edge}; the zero-defect endpoint is included
by the zero-radius cap conclusion.  The retained estimate
$H_d^{1/(d-1)}\leq144$, together with $d+1\leq4d/3$, gives
\eqref{eq:inverse-actual-edge-coarse}, since
$192(4+2\sqrt2)=768+384\sqrt2$.
For completeness, the numerical bound on $H_d$ follows from
$H_3=8640<144^2$ and
$H_{d+1}/H_d\leq25/6<144$ for $d\geq3$, by induction.
If $e(K)>H_d^{-1}$, then $H_d^{1/(d-1)}e(K)^{1/(d-1)}>1$,
so the universal $E(K,S)\leq d+1$ proves
\eqref{eq:inverse-actual-edge} directly; the coarser bound follows.
\end{proof}

\begin{remark}[The compatibility quantifier]
\label{rem:inverse-compatible-P}
The hypothesis refers to a $P$ produced by the actual-defect chain.
The existence of some unrelated enclosing simplex with edge-supported
extreme points is insufficient.  No result here shows that the
chain's $P$ can always be chosen with this property.  If only one
prescribed $S$ is to be retained, it is enough that this $S$ be
edge-supported relative to an admissible $P$, by
Theorem~\ref{thm:inverse-edge-simplex}.  The body-level hypothesis is
what gives the conclusion simultaneously for every prescribed
maximum simplex without inspecting each support.  The exact
coefficient in \eqref{eq:inverse-actual-edge} is asymptotic to
$(4+2\sqrt2)d$, since $H_d^{1/(d-1)}\to1$.  These are conditional
corollaries of the original invariant and the retained upstream
chain, not an unconditional theorem for all convex bodies.
\end{remark}

\subsection{Order three and an exact golden-ratio obstruction}

\begin{theorem}[An order-three two-minor inequality]
\label{thm:inverse-order-three}
If $n=3$, then
\begin{equation}
 1-\min_{i,j}B_{ij}\leq M_2(B).
 \label{eq:inverse-order-three}
\end{equation}
Consequently, if $M_1(B)\leq c$ and $M_2(B)\leq c^2$, then
\begin{equation}
 a(B)\leq\min\{c,c^2-1\}\leq\varphi^2(c-1),
 \qquad \varphi=\frac{1+\sqrt5}{2}.
 \label{eq:inverse-order-three-linear}
\end{equation}
The coefficient $\varphi^2$ is optimal.  The same conclusions hold
for direct sums of blocks of orders at most three, after independent
row and column permutations.
\end{theorem}

\begin{proof}
Permutations preserve stochasticity and the quantities in
\eqref{eq:inverse-order-three}.  First arrange
$D=\det W>0$.  Independent even row and column permutations are
transitive on the three indices, so any chosen entry of $B$ can be
placed at position $(3,1)$ while preserving $D>0$.  Write
\[
 W=\begin{pmatrix}
 u&v&z\\r&s&t\\1-u-r&1-v-s&1-z-t
 \end{pmatrix}.
\]
If $C_{13}$ is the cofactor at row $1$, column $3$, direct
expansion gives
\begin{align*}
 D-C_{13}
 &=(s-r)(1-z)+t(v-u)\\
 &=(1-z-t)(s-r)+t[(v+s)-(u+r)]+z\cdot0.
\end{align*}
This is a convex combination, with the third-column weights
$1-z-t,t,z$, of the values
$s-r$, $W_{31}-W_{32}$, and $0$.  Each is at most
$w_*:=\max_{i,j}W_{ij}$.  Since $B_{31}=C_{13}/D$,
\[
 1-B_{31}\leq\frac{w_*}{D}.
\]
Jacobi's complementary-minor identity gives $M_2(B)=w_*/D$.
Applying the argument to every entry proves
\eqref{eq:inverse-order-three}.  Thus $a(B)\leq c^2-1$, and the
entry bound gives $a(B)\leq c$.  For $c>1$, the upper envelope
of $\min\{c,c^2-1\}/(c-1)$ has maximum $\varphi^2$ at the crossing
$c=\varphi$; $c=1$ follows directly.

A one-by-one block is $[1]$.  In order two, every
column-stochastic inverse can be written
\[
 \begin{pmatrix}x&y\\1-x&1-y\end{pmatrix}.
\]
Each negative entry has a positive entry in its column larger in
absolute value by one.  Hence $a(B)\leq M_1(B)-1\leq c-1$ in
that order.  Applying the appropriate bound in each block proves
the direct-sum assertion.  Optimality follows from the next exact
example.
\end{proof}

\begin{proposition}[An all-order golden-ratio extremizer]
\label{prop:inverse-golden-example}
Let $p=\varphi^{-1}$ and set
\[
 W=\begin{pmatrix}
 0&p&1/2\\p&0&1/2\\1-p&1-p&0
 \end{pmatrix}.
\]
Then $W$ is nonnegative and column-stochastic,
$\det W=p(1-p)=\varphi^{-3}$, and
\[
 B=W^{-1}=\begin{pmatrix}
 -\varphi/2&\varphi/2&\varphi^2/2\\
 \varphi/2&-\varphi/2&\varphi^2/2\\
 1&1&-\varphi
 \end{pmatrix}.
\]
Its all-order root bound is exactly $c=\varphi$, and
\begin{equation}
 M_1(B)=\varphi,\quad M_2(B)=\varphi^2,\quad
 M_3(B)=\varphi^3,\quad
 a(B)=\varphi=c^2-1=\varphi^2(c-1).
 \label{eq:inverse-golden-equality}
\end{equation}
Thus any universal estimate $a(B)\leq C(c-1)$ valid for all
$c\geq1$ must have $C\geq\varphi^2$.
\end{proposition}

\begin{proof}
The stochasticity and displayed inverse follow by direct
multiplication using $\varphi^2=\varphi+1$.  The largest absolute
entry of $B$ is $\varphi$.  Jacobi gives
$M_2(B)=p/\det W=\varphi^2$, and
$M_3(B)=1/\det W=\varphi^3$.  The minimum entry is $-\varphi$,
which yields all equalities in \eqref{eq:inverse-golden-equality}.
\end{proof}

This is an exact obstruction to a smaller global constant.  It is
not a counterexample in the limit $c\downarrow1$.  The order-three
inequality \eqref{eq:inverse-order-three} is dimension-specific;
the corresponding proposed estimate with
$\max\{M_1(B),M_2(B)\}$ is false in order five.  In particular, an
unrestricted extension based only on first- and second-order minors
would require a new argument.

\subsection{Positive definiteness and a local permutation basin}

\begin{theorem}[Symmetric positive-definite case]
\label{thm:inverse-positive-definite}
If $W$ is symmetric and positive definite, then $M_1(B)\leq c$
alone implies
\begin{equation}
 a(B)\leq c-1.
 \label{eq:inverse-PSD}
\end{equation}
The coefficient one is sharp even under all-order minor bounds.
In the setting of Lemma~\ref{lem:inverse-geometric-reduction}, a
symmetric positive-definite prescribed $W$ gives the all-scale bound
\begin{equation}
 E(K,S)\leq4n\rho.
 \label{eq:inverse-PSD-geometric}
\end{equation}
\end{theorem}

\begin{proof}
A nonnegative column-stochastic matrix has spectral radius at most
one.  Symmetry and positive definiteness put all eigenvalues of $W$
in $(0,1]$.  Therefore $H=B-\Id_n$ is positive semidefinite and
\[
 0\leq H_{ii}=B_{ii}-1\leq c-1.
\]
Every two-coordinate positive-semidefinite restriction gives, for
$i\ne j$,
\[
 |B_{ij}|=|H_{ij}|\leq\sqrt{H_{ii}H_{jj}}\leq c-1.
\]
The diagonal entries are at least one, so this proves
\eqref{eq:inverse-PSD}.

For sharpness take
$W=\left(\begin{smallmatrix}1-s&s\\s&1-s\end{smallmatrix}\right)$
with $0<s<1/2$.  Its inverse has negative magnitude
$a=s/(1-2s)$, diagonal $1+a$, and determinant $1+2a$.
Since $(1+a)^2\geq1+2a$, its all-order root bound is $c=1+a$.
Thus equality holds in \eqref{eq:inverse-PSD}.

For the geometric claim, if $0\leq\rho<1/2$ then
\eqref{eq:inverse-all-body-excess} gives
$E(K,S)/n\leq\min\{1,c-1\}$ with $c=(1-2\rho)^{-1}$.
For $\rho\leq1/4$,
$c-1=2\rho/(1-2\rho)\leq4\rho$; for $\rho\geq1/4$ the
universal bound $E(K,S)/n\leq1\leq4\rho$ applies, also beyond
the Neumann range.
\end{proof}

The symmetry and positive-definiteness assumption is
coordinate-dependent.  It must be verified and does not follow
from positivity of the entries of $W$.

\begin{lemma}[A local permutation basin]
\label{lem:inverse-local-basin}
Suppose independent row and column permutations put $W$ in a form
with
\[
 s=\max_j(1-W_{jj})<\frac14.
\]
Then $M_1(B)\leq c$ implies
\begin{equation}
 a(B)\leq\frac{c-1}{1-4s}.
 \label{eq:inverse-local-basin}
\end{equation}
In particular, $s\leq1/8$ gives $a(B)\leq2(c-1)$.
\end{lemma}

\begin{proof}
Put $H=\Id_n-W$.  Its induced column norm is
$\|H\|_1=2s$, and its maximum entry modulus is at most $s$.
For $k\geq1$,
\[
 \max_{i,j}|(H^k)_{ij}|\leq s(2s)^{k-1}.
\]
The Neumann series therefore gives $a(B)\leq s/(1-2s)$.
At an index $j$ with $H_{jj}=s$, it also gives
\[
 B_{jj}\geq1+s-\frac{2s^2}{1-2s}
 =1+\frac{s(1-4s)}{1-2s}.
\]
Since $B_{jj}\leq c$, combining these inequalities proves
\eqref{eq:inverse-local-basin}.  The case $s=0$ is immediate.
\end{proof}

This lemma does not assert the existence of a suitable permutation.
A dimension-free argument placing arbitrary instances in such a
basin would itself constitute a substantive missing step.

\subsection{Additional all-order reductions}

The following results preserve useful information beyond the
support-two theorem.  Here the hypothesis is again
$M_k(B)\leq c^k$ for every relevant order.

\begin{proposition}[Localized minor completion and occupancy]
\label{prop:inverse-occupancy}
Let $R,S\subset[n]$ be disjoint row-label sets, with
$|R|=k$ and $|S|=\ell$, and let $J\subset[n]$ have size $k$.
For each $s\in S$, independently draw $Y_s$ from the distribution
$W_{:,s}$.  Then
\begin{equation}
 |\det B_{R,J}|
 \leq c^{k+\ell}\mathbb P\bigl(
     Y_s\text{ are pairwise distinct and all outside }J\bigr).
 \label{eq:inverse-occupancy}
\end{equation}
In particular, for any $B_{ij}$ and $S\subset[n]\setminus\{i\}$,
\begin{align}
 |B_{ij}|&\leq c^{|S|+1}
  \mathbb P\bigl(Y_s\text{ are distinct and all different from }j\bigr)
  \nonumber\\
 &\leq c^{|S|+1}\prod_{s\in S}(1-W_{j,s}).
 \label{eq:inverse-occupancy-entry}
\end{align}
With $k=0$ and the empty determinant equal to one,
\begin{equation}
 1\leq c^{|S|}\mathbb P(Y_s\text{ are pairwise distinct}).
 \label{eq:inverse-injection}
\end{equation}
\end{proposition}

\begin{proof}
Order $R\cup S$ with $R$ first.  Form the $(k+\ell)$-square
matrix whose first $k$ columns are $B_{R\cup S,J}$ and whose
last $\ell$ columns are the coordinate vectors $e_s$ restricted to
$R\cup S$.  Its determinant is $\det B_{R,J}$ up to a fixed
ordering sign.  Substitute
\[
 e_s=\sum_h W_{h,s}B_{:,h}
\]
and expand multilinearly in the last $\ell$ columns.  A term
vanishes if its selected indices repeat or if one belongs to $J$.
Every remaining determinant is a $(k+\ell)$-minor of $B$, bounded
in magnitude by $c^{k+\ell}$.  The sum of the corresponding
nonnegative product weights is exactly the probability in
\eqref{eq:inverse-occupancy}.  Taking $k=1$ gives the first entry
bound.  Dropping the collision-free condition and using the
independent draws gives the product bound.  Taking $k=0$ gives
\eqref{eq:inverse-injection}.
\end{proof}

The collision-free term retains both the fixed minor and its
forbidden column set and can be stronger than the displayed product
bound.

\begin{proposition}[Regularization of a uniformly almost nonnegative inverse]
\label{prop:inverse-regularization}
Suppose $A$ is invertible, every square minor of $A$ has absolute
value at most one, $\ones^{\mathsf T}A=\ones^{\mathsf T}$, and
$V=A^{-1}$ satisfies
$V_{ij}\geq-\eta/n$ for some $\eta\geq0$.  With
$J=\ones\ones^{\mathsf T}$, define
\[
 W=\frac{V+(\eta/n)J}{1+\eta},\qquad
 Q=(1+\eta)\Id_n-\frac\eta n J,\qquad B=W^{-1}=AQ.
\]
Then $W$ is nonnegative and column-stochastic,
\begin{equation}
 |B_{ij}-A_{ij}|\leq2\eta,
 \qquad M_k(B)\leq(1+2\eta)^k\quad(1\leq k\leq n).
 \label{eq:inverse-regularization}
\end{equation}
Thus a universal bound $a(B)\leq C(c-1)$ in the stochastic
inverse-minor problem would imply
$-\min A\leq(2C+2)\eta$ under these hypotheses.
\end{proposition}

\begin{proof}
The assumed lower bound gives $W\geq0$, and
$\ones^{\mathsf T}V=\ones^{\mathsf T}$ gives its column sums.
Using the same column-sum identity for $A$ verifies $B=AQ$.
Entrywise,
\[
 B_{ij}-A_{ij}
 =\eta\left(A_{ij}-\frac1n\sum_h A_{ih}\right).
\]
The order-one minor bound on $A$ proves the first part of
\eqref{eq:inverse-regularization}.

For a fixed $k$-element column set, the only nonzero $k$-minors
of $Q$ have the same row set or exchange one index.  The principal
minor is
\[
 (1+\eta)^{k-1}(1+\eta-k\eta/n),
\]
and each of the $k(n-k)$ one-index-exchange minors has magnitude
$(\eta/n)(1+\eta)^{k-1}$.  Cauchy--Binet therefore gives
\begin{align*}
 |\det B_{I,J}|&\leq
 (1+\eta)^{k-1}
 \left[1+\eta+\frac{\eta k(n-k-1)}n\right]\\
 &\leq(1+\eta)^{k-1}[1+(k+1)\eta]
 \leq(1+2\eta)^k.
\end{align*}
For the last inequality, Bernoulli's inequality gives
\[
 \frac{(1+2\eta)^k}{(1+\eta)^{k-1}}
 =(1+\eta)\left(1+\frac\eta{1+\eta}\right)^k
 \geq1+(k+1)\eta.
\]
Applying the hypothetical stochastic bound with $c=1+2\eta$
and then the entrywise difference estimate gives
$-\min A\leq2C\eta+2\eta$, as asserted.
\end{proof}

This is a one-way reduction.  It does not establish that an
arbitrary stochastic instance admits the stronger regularized
representation in the opposite direction.

\begin{lemma}[Conditional stability from a dominant matching]
\label{lem:inverse-dominant-matching}
If, after a row permutation, $W_{ii}\geq1-\varepsilon$ for every
$i$, with $0\leq\varepsilon<1/2$, then
\begin{equation}
 \|W^{-1}-\Id_n\|_1\leq
 \frac{2\varepsilon}{1-2\varepsilon}.
 \label{eq:inverse-dominant-matching}
\end{equation}
In particular the same upper bound holds for the negative entry
magnitudes in the permuted inverse.
\end{lemma}

\begin{proof}
Nonnegativity and column stochasticity give
$\|\Id_n-W\|_1\leq2\varepsilon$.  The Neumann series yields
\[
 \|W^{-1}-\Id_n\|_1
 \leq\frac{\|\Id_n-W\|_1}{1-\|\Id_n-W\|_1}
 \leq\frac{2\varepsilon}{1-2\varepsilon}.
\]
Every negative entry is bounded by this norm, proving the claim.
\end{proof}

A dimension-free dominant-matching theorem of the form
$W_{i,\sigma(i)}\geq1-O(c-1)$ would settle the target through
this lemma.  No such matching theorem is established here.

\subsection{Limitations of the structural results}

The support-two proof uses the unicyclic graph decomposition
essentially.  A column supported on three or more rows does not
reduce to the same path recursion, and no reduction of arbitrary
$W$ to this class is supplied.  The all-order conjecture
$a(B)\leq c^2-1$ remains open in general; its analogue based only
on first- and second-order bounds is false in order five.
Neither an unrestricted dimension-free inverse-minor theorem nor
a counterexample to it follows from this section.  The special
central-$Q_t$ formulation from Section~\ref{sec:projection} and
the unrestricted prescribed-maximum linear dimension bound also
remain open.  In particular, a matrix construction by itself is
not an actual-defect lower-bound body: the all-direction projection
and enclosing-simplex requirements must still be verified.

\appendix
\section{A finite determinant assignment theorem}
\label{sec:assignment}
We record a finite version of the retained volume-weighted assignment
argument, with its proof. For a real number $t$, write $t_+=\max(t,0)$.
All probability measures in this section have finite support, so no
integrability or limiting-law issue is present.

\begin{lemma}[Seminorm assignment]\label{lem:assignment}
Let $\mu$ be a centered probability measure with finite support
$\mathcal A\subset\R^d$ whose affine span is $\R^d$. Let $N$ be a seminorm
with $N(x)\le R_0$ on $\mathcal A$, where $R_0\ge0$. For independent
samples with law $\mu$, put
\[
 A=\E\abs{\det(X_1,\ldots,X_d)},\qquad
 B=\E\abs{\det((X_0,1),\ldots,(X_d,1))}.
\]
Then $B>0$, $D:=B-A\ge0$, and there are affinely independent anchors
$w_0,\ldots,w_d\in\mathcal A$ and an assignment $I:\mathcal A\to\{0,\ldots,d\}$
with
\begin{equation}\label{eq:assignment}
 h:=\E N(X-w_{I(X)})\le R_0(d+1)(d+2)\frac{D}{B}.
\end{equation}
One can choose $I(x)$ to be the least index of a largest barycentric
coordinate of $x$ in the anchor simplex.
\end{lemma}
\begin{proof}
For an ordered base $\mathbf x=(x_1,\ldots,x_d)$ define
\[
 F_{\mathbf x}(y)=\det((x_1,1),\ldots,(x_d,1),(y,1)),\quad
 P_{\mathbf x}=\E(F_{\mathbf x}(X))_+,\quad
 M_{\mathbf x}=\E(-F_{\mathbf x}(X))_+.
\]
Centering and affinity in the last column give
$\abs{\E F_{\mathbf x}(X)}=\abs{\det(x_1,\ldots,x_d)}$. Hence
\begin{equation}\label{eq:Dsign}
 D=2\E_{\mathbf X}\min(P_{\mathbf X},M_{\mathbf X})\ge0.
\end{equation}
For two test points set
\[
 \psi_{\mathbf x}(y,z)=
 \min((F_{\mathbf x}(y))_+,(-F_{\mathbf x}(z))_+)
 +\min((-F_{\mathbf x}(y))_+,(F_{\mathbf x}(z))_+).
\]
Each conditional expected minimum is at most
$\min(P_{\mathbf x},M_{\mathbf x})$. Therefore
\begin{equation}\label{eq:psibudget}
 \E_{\mathbf X,Y,Z}\psi_{\mathbf X}(Y,Z)\le D.
\end{equation}
These statements include singular bases.

For an anchor tuple $W=(w_0,\ldots,w_d)$, put
$v(W)=\abs{\det((w_0,1),\ldots,(w_d,1))}$. Form $G(W,x)$ by adding the
following witnesses, with any fixed order on each base:
\begin{enumerate}[label=(\roman*),leftmargin=2em]
\item for each $i$, the base is the other $d$ anchors and the test points
are $w_i,x$;
\item for each $i<j$, the base is $x$ and the other $d-1$ anchors, and the
test points are $w_i,w_j$.
\end{enumerate}
Reversing base orientation does not change $\psi$. There are
$M_d=(d+1)(d+2)/2$ summands. Each uses $d+2$ independent samples when
$W$ and $x$ are sampled, so, with $J(W)=\E_XG(W,X)$,
\[
 \E_WJ(W)\le M_dD,\qquad \E_Wv(W)=B>0.
\]
Positivity of $B$ follows from an affinely independent tuple in the support.
Weighting anchor tuples by $v(W)/B$ gives an anchor tuple with $v(W)>0$ and
\begin{equation}\label{eq:weightedchoice}
 \frac{J(W)}{v(W)}\le M_d\frac{D}{B}.
\end{equation}
Indeed its weighted mean is
$\E[J(W)\mathbf1_{\{v(W)>0\}}]/B\le M_dD/B$.
No assertion that $J$ vanishes at singular tuples is needed.

Fix this tuple, and write $\alpha_i(x)$ for its barycentric coordinates.
Determinant affinity gives, for witnesses of type (i),
$v(W)\min(1,(-\alpha_i)_+)$; those of type (ii) equal
\[
 v(W)\bigl[\min((\alpha_i)_+,(\alpha_j)_+)
           +\min((-\alpha_i)_+,(-\alpha_j)_+)\bigr].
\]
It follows that $G(W,x)\ge v(W)\phi(x)$, where
\[
 \phi(x)=\sum_i\min(1,(-\alpha_i(x))_+)
       +\sum_{i<j}\min((\alpha_i(x))_+,(\alpha_j(x))_+).
\]
Let $j=I(x)$ be a largest coordinate. Since $\sum_i\alpha_i=1$, we have
$\alpha_j>0$. Put $N_x=\sum_i(-\alpha_i)_+$ and
$U_x=\sum_{i\ne j}(\alpha_i)_+$. The support bounds yield both
\[
 N(x-w_j)\le2R_0(N_x+U_x),\qquad N(x-w_j)\le2R_0.
\]
The first follows by expanding $x-w_j=\sum_{i\ne j}\alpha_i(w_i-w_j)$;
the second uses $x,w_j\in\mathcal A$. The pairs containing the largest
coordinate give $U_x\le\sum_{i<l}\min((\alpha_i)_+,(\alpha_l)_+)$.
Also $\min(1,N_x)\le\sum_i\min(1,(-\alpha_i)_+)$. Thus
\[
 N(x-w_{I(x)})\le2R_0\min(1,N_x+U_x)
 \le2R_0\phi(x)\le2R_0\frac{G(W,x)}{v(W)}.
\]
Integrating and using \eqref{eq:weightedchoice} proves \eqref{eq:assignment}.
The reasoning includes $D=0$ without division by $D$.
\end{proof}

\section{From the actual facet law to an enclosing simplex}
\label{sec:enclosure}
\subsection{Raw facet identities and the original defect}\label{prop:actual-moments}
For a full-dimensional $p$-polytope $L$, let $u_i=s_i n_i$ be its
facet area-normal vectors and $b_i=s_i h_L(n_i)$, where $n_i$ is the
outward unit normal and $s_i>0$ is the facet area. Set
\begin{equation}\label{eq:rawsums}
 \mathcal H(L)=\sum_{|J|=p}\abs{\det(u_j:j\in J)},\qquad
 \mathcal L(L)=\sum_{|J|=p+1}\abs{\det((u_j,b_j):j\in J)}.
\end{equation}
The order chosen within a subset is immaterial. Translating $L$ replaces
$b_i$ by $b_i-\ip{u_i}{z}$ and leaves these sums unchanged.

Cauchy's formula gives $b_L(z)=\frac12\sum_i\abs{\ip{z}{u_i}}$.
Hence $\Pi L=\sum_i[-u_i/2,u_i/2]$, and the usual volume formula for a
zonotope gives $|\Pi L|=\mathcal H(L)$. Direct facet parametrization of the
unit-height pyramid gives its side pairs and bottom pair as
\begin{equation}\label{eq:pyrfacets}
 \left(\frac{(u_i,b_i)}p,\frac{b_i}p\right),\qquad ((0,-|L|),0),
\end{equation}
respectively. One way to check the factor $1/p$ is to integrate the
$(p-1)$-dimensional slices of each side face: its oriented area vector is
$(u_i,b_i)\int_0^1(1-t)^{p-1}dt=(u_i,b_i)/p$.
Expanding the corresponding minors gives
\begin{align}
 \mathcal H(\pyr L)&=p^{-(p+1)}\mathcal L(L)
                     +|L|p^{-p}\mathcal H(L),\label{eq:pyrH}\\
 \mathcal L(\pyr L)&=|L|p^{-(p+1)}\mathcal L(L),
 &|\pyr L|&=\frac{|L|}{p+1}.\label{eq:pyrL}
\end{align}
For the second identity, an all-side lifted minor is zero because its
last two coordinate rows coincide; minors using the bottom give the
stated expression. Substituting \eqref{eq:pyrH} into the definition of $a$
therefore proves the exact identity
\begin{equation}\label{eq:a_raw}
 a(L)=\frac{\mathcal L(L)}{p|L|\mathcal H(L)}.
\end{equation}
These computations concern the original projection bodies, not a substitute
invariant. The standard projection covariance
$\Pi(AL)=|\det A|A^{-T}\Pi L$ shows that $R$, and hence $a$ and $e$,
are affine invariant; affine maps of a base extend to its pyramid.
The quantity $E$ is affine invariant as well.

\subsection{Normalization and the finite cone law}
Affinely normalize the prescribed maximum $S$ so that its own centroid is
zero and its image contains the closed unit ball. Continue to write $K,S$
for the images. With $n=d+1$, maximum-volume replacement gives
\begin{equation}\label{eq:replacement}
 |\alpha_i^S(x)|\le1\quad(x\in K).
\end{equation}
Indeed replacing vertex $i$ by $x$ multiplies the absolute volume by
$|\alpha_i^S(x)|$. Since $\sum_i\alpha_i^S(x)=1$ and the centroid is zero,
$-x/d$ has $S$-coordinates $(1-\alpha_i^S(x))/d\ge0$. Thus
\begin{equation}\label{eq:asymmetry}
 -K\subset dS\subset dK,\qquad E(K,S)\le n.
\end{equation}
The second conclusion uses the exact centered-dilation description
\begin{equation}\label{eq:dilationcoords}
 g_S+(1+\tau)(S-g_S)=\{x:\alpha_i^S(x)\ge-\tau/n\text{ for all }i\}.
\end{equation}

For the actual irredundant facets $F$ of $K$, use the notation
$n_F,s_F,h_F=h_K(n_F)$ and put
\begin{equation}\label{eq:law}
 X_F=\frac{n_F}{h_F},\qquad p_F=\frac{s_Fh_F}{d|K|},\qquad
 \mu=\sum_F p_F\delta_{X_F}.
\end{equation}
Here $h_F\ge1$ because the unit ball is in $K$. The cone-volume identity
$\sum_Fs_Fh_F=d|K|$ and the surface-area balance $\sum_Fs_Fn_F=0$ show
that $\mu$ is a centered probability law. Its atoms are precisely actual
facet normals divided by their positive support numbers, with
$\norm{X_F}\le1$ and $h_K(X_F)=1$.
Boundedness of $K$ forces these atoms to span $\R^d$; centeredness with
positive weights then gives full affine span. In particular the horizontal
and lifted moments $A,B$ of Section~\ref{sec:assignment} are positive.

Cauchy's formula and centeredness give the exact brightness identities
\begin{equation}\label{eq:brightness}
 m_K(z):=\E(\ip{z}{X})_+=\frac{b_K(z)}{d|K|},\qquad
 f_K(z):=\E|\ip zX|=2m_K(z).
\end{equation}
Expanding the iid determinants by distinct facet subsets gives
\[
 A=\frac{d!\mathcal H(K)}{(d|K|)^d},\qquad
 B=\frac{(d+1)!\mathcal L(K)}{(d|K|)^{d+1}}.
\]
Together with \eqref{eq:a_raw}, this proves
\begin{equation}\label{eq:momentdefect}
 a(K)=\frac{B}{nA},\qquad
 \frac{B-A}{B}=\frac{ne(K)}{1+ne(K)}.
\end{equation}
Jensen's inequality, or \eqref{eq:Dsign}, proves $B\ge A$ and hence $e(K)\ge0$.

Let $N=h_{K-K}$, a continuous seminorm. At every atom,
\[
 N(X_F)=h_K(X_F)+h_K(-X_F)\le1+d=n
\]
by \eqref{eq:asymmetry}. Lemma~\ref{lem:assignment}, with $R_0=n$, gives
$d+1$ affinely independent \emph{actual atoms} $w_i$ and an assignment with
\begin{equation}\label{eq:hbudget}
 h=\E N(X-w_{I(X)})\le
 \frac{Ce(K)}{1+ne(K)},\qquad C=n^3(n+1).
\end{equation}
Thus every selected $w_i$ is an actual facet normal, a fact that is stronger
than merely knowing that an anchor belongs to an unspecified support law.

\subsection{Width and same-anchor enclosure}
For every $z,y\in\R^d$ one has
\begin{equation}\label{eq:width}
 |\ip zy|\le d\,m_K(z)N(y).
\end{equation}
To see this, first take unit $v$ and $c=\ip yv\ne0$. In coordinates
$x=q+tv$ with $q\in v^\perp$, the determinant-one shear
$(q,t)\mapsto(q,t+\ip yq/c)$ places $K$ in a cylinder with base the
projection of $K$ and height $N(y)/|c|$. Thus
$|c||K|\le N(y)\pi_K(v)$. Scaling $v$ to $z$ and using
\eqref{eq:brightness} proves \eqref{eq:width}; zero cases are immediate.

Let $p_i=\mu\{I=i\}$. Positive-part Lipschitz continuity gives
\begin{equation}\label{eq:positivecomp}
 \sum_i p_i(\ip z{w_i})_+\ge(1-dh)m_K(z).
\end{equation}
Assume $h<1/(2d)$. On the unit sphere, $m_K$ has a positive lower bound,
since $K$ contains the unit ball. Hence for every unit $z$ some anchor has
positive pairing with $z$, uniformly bounded away from zero. Support
separation puts zero in the interior of $\conv(w_0,\ldots,w_d)$.
Its polar
\begin{equation}\label{eq:outer}
 P=\{x:\ip{w_i}x\le1\text{ for }0\le i\le d\}
\end{equation}
is therefore a genuine $d$-simplex containing $K$. Every defining
inequality is an actual $K$-facet inequality. The corresponding full facet
of $K$ lies in a facet of $P$, providing the actual-facet incidence used in Theorem~\ref{thm:few-facet}.

For $z\in P$, the left side of \eqref{eq:positivecomp} is at most one;
therefore $m_K(z)\le2$. Equation~\eqref{eq:width} and support separation
then give the intrinsic containment
\begin{equation}\label{eq:intrinsiccontain}
 P\subset2d(K-K).
\end{equation}

\subsection{Probability correction and relative projection error}
Let $\lambda_i>0$ be the barycentric coordinates of zero in the anchor
simplex. Let $q_i$ be the vertex of $P$ opposite the facet with normal $w_i$.
The anchor-simplex coordinates satisfy
\begin{equation}\label{eq:polarcoords}
 \alpha_i^w(x)=\lambda_i(1-\ip{q_i}x).
\end{equation}
Both sides agree at each anchor. The actual cone probabilities of $P$
satisfy the same mass and mean equations as $\lambda_i$ and hence equal
$\lambda_i$. Put $c=\sum_i p_iw_i$. Since $\E X=0$,
\eqref{eq:intrinsiccontain} gives $|\ip{q_i}c|\le2dh$. Evaluating
\eqref{eq:polarcoords} at $c$ yields
\[
 p_i=\lambda_i(1-\ip{q_i}c)\ge(1-2dh)\lambda_i.
\]
Also \eqref{eq:width} and the absolute-value triangle inequality give
\[
 \sum_i p_i|\ip z{w_i}|\le f_K(z)+d m_K(z)h
 =(1+dh/2)f_K(z).
\]
Write $\theta=2dh<1$. The exact brightness formulas for $P$ and $K$ imply
\begin{equation}\label{eq:brightnesscomp}
 (1-\theta)\frac{\pi_P(u)}{|P|}
 \le(1+\theta/4)\frac{\pi_K(u)}{|K|}\qquad(\norm u=1).
\end{equation}

The finite first mixed-volume formula and Minkowski's first inequality give
\begin{equation}\label{eq:volumeroot}
 \left(\frac{|P|}{|K|}\right)^{1/d}
 \le\frac{V(K[d-1],P)}{|K|}
 =\E h_P(X).
\end{equation}
Indeed the last expression is
$\sum_Fs_Fh_P(n_F)/(d|K|)$. Since $h_P(w_i)=1$, subadditivity and
\eqref{eq:intrinsiccontain} bound it by $1+2dh=1+\theta$.
Combining \eqref{eq:brightnesscomp} and \eqref{eq:volumeroot},
\[
 \frac{\pi_K(u)}{\pi_P(u)}
 \ge\frac{1-\theta}{(1+\theta)^d(1+\theta/4)}.
\]
For factors in $[0,1]$, the deficit of their product is at most the sum of
their deficits. Thus the missing fraction above is at most
\begin{equation}\label{eq:scalar}
 \theta+d\theta+\theta/4
 =2d(d+5/4)h\le3d^2h\qquad(d\ge3).
\end{equation}
Here $1-(1+\theta)^{-d}\le d\theta$ follows from
$1-t^d\le d(1-t)$ for $0\le t\le1$.

\begin{proposition}[Finite actual-facet enclosure]\label{prop:enclosure}
For $d\ge3$, if $0\le e(K)\le H_d^{-1}$, the normalized $K$ is contained
in a simplex $P$ whose facet hyperplanes are actual $K$-facet hyperplanes,
and
\begin{equation}\label{eq:finalprojection}
 \pi_P(u)-\pi_K(u)\le H_de(K)\pi_P(u)\qquad(\norm u=1).
\end{equation}
\end{proposition}
\begin{proof}
Since $H_d=3d^2C$, \eqref{eq:hbudget} gives
$h\le1/(3d^2)<1/(2d)$ even at $e(K)=H_d^{-1}$. All strict enclosure gates
hold. Equations~\eqref{eq:scalar} and \eqref{eq:hbudget} give
$3d^2h\le H_de(K)$.
\end{proof}
Every step used the same finite actual facet law. No compact approximation,
subsequence, or limiting-law identification is needed for this proposition.

\section{Proof coverage and verification boundaries}\label{app:review}
\subsection{Mathematical provenance}
The manuscript organizes the research series dated 8 October 2026,
whose frozen source archive has SHA-256
\begin{center}\small\ttfamily
631a80b594fb9a597fef15e1c76870956ec\\b768fbbb14f3b7953e1809970296e.
\end{center}
The companion source map identifies each statement with its mathematical
source and a version-bound review. The 55 source-manifest payload files
are preserved byte for byte; the source archive also contains its manifest.
This manuscript is a new exposition and does not overwrite any frozen
source or earlier manuscript.

Sections~\ref{sec:clipping}--\ref{sec:few} reproduce the complete four-note
few-facet series: rooted-forest determinants and equality, arbitrary
nonvertex maxima, original-centroid excess, arbitrary one-corner bodies,
actual signed-cofactor identities, and the full dimensional squeeze.
The lower-body facet certificate comes from the retained realized-body
obstruction note. The exact truncation identity is included rather than
using an unidentified geometric surrogate for the defect.

Section~\ref{sec:projection} includes the elementary and
Brunn--Minkowski core proofs, the finite-spectrum refinement and its
uniform branch, the beta/binomial face hierarchy, and the necessary
countermodel tests. Section~\ref{sec:inverse} includes weighted cycles,
attached trees, the replacement-minor reduction, averaging for every
prescribed maximum, the compatible actual-defect consequence, and the
additional order-three, positive-definite, and local-basin statements.
The latter finite structures do not establish unrestricted retention.

Appendices~\ref{sec:assignment} and \ref{sec:enclosure} reproduce the
finite determinant-assignment and actual-facet enclosure proofs from the
existing vertex-excess manuscript exposition, adapted only to this
paper's references. They expand the exact finite-law interface identified
in the frozen series. In particular, facet incidence is proved by selection
from actual irredundant facets; it is not claimed as an unproved consequence
of a compact-limit formal theorem. Classical mixed-volume and projection
formulas are explicit standard inputs, not newly proved stability results.

\subsection{The traditional review record}
The four-note few-facet series received a separate independent traditional
model review with a PASS disposition bound to the frozen versions.
That review checked all four proofs, the arbitrary prescribed-$S$
quantifier, the original invariant, the finite-facet interface, and the
realized lower examples. It inspected the retained interface without
re-proving the entire upstream actual-body theory. Its only source edit
was the resolved replacement of ``span'' by ``form a basis'' in the
rooted-forest rank argument. This manuscript follows the corrected version.

The full-projection series received the parent research worker's
traditional proof and integration review. The uniform $\rho\le1/3$
branch was contributed by that parent and cross-checked by the originating
projection worker. Thus the parent review is not independent review of
its own authored refinement, and the entire series is not presented as
having a uniformly independent authorship-and-review arrangement.

The inverse-minor series received a parent traditional reading and a
separate theoretical worker's proof work and focused integration audit.
The latter's final scope includes the all-body coordinate control,
replacement-minor reduction, every-maximizer averaging, compact
one-skeleton reduction, conditional actual-defect consequence, and the
order-three and golden-ratio formulas. The corrected hypothesis uses
\emph{every extreme point} on the admissible simplex's one-skeleton;
an earlier ambiguous wording is not retained. The separate focused audit
treats the upstream actual enclosure chain as an input. Because the
worker also contributed proof material, it is not described here as a
wholly external independent review of every matrix proof.

All of these are traditional model-assisted mathematical checks, not
professional peer review. They bind source versions. They do not by
themselves certify the typesetting, cross-references, or transcription of
this new manuscript. Its build, page-by-page visual inspection, proof
coverage, and source hashes are recorded in the companion candidate QA
report; final independent copy review remains a separate gate.

\subsection{Exact finite checks}
The source review record contains the following reproducible checks.
\begin{itemize}[leftmargin=1.5em]
\item The clipping checker covers 12 rational examples and all 12,769
vertex tuples. An independent Bareiss implementation verifies their full
forest product and equality classification: 8,050 singular tuples,
4,719 rooted forests, 64 maxima, and 4,655 strict nonmaximum forests.
\item The author cofactor checker covers 60 balanced rational matrices
and 32 actual-truncation comparisons. A separate integer checker covers
210 augmented matrices and 4,830 omitted-pair identities. General matrix
samples test algebra, not geometric realizability as facet data.
\item Projection checks cover 189 exact beta/binomial identities and
780 powered rational comparisons for the uniform spectrum branch in
$2\le m\le40$. The all-dimensional conclusion uses the written scalar
proof, not these finite samples.
\item Inverse-minor replay covers 140 order-three and 287 support-two
rational matrices, including 149 near-endpoint cases, the symbolic
golden-ratio identities, averaging and cap constants, and the known gaps
in two previously proposed false unrestricted two-minor bounds.
\end{itemize}
The author's four certificate outputs were replayed byte identically in
the frozen series. Exact finite tests support arithmetic and implementation
consistency; they do not replace the universal proofs or certify novelty.

\subsection{Formal status and the unresolved question}
The earlier unrestricted quadratic-dimensional result has a retained
strict Lean compilation and same-implementation kernel-replay record.
No new Lean compilation, kernel replay, or formal bridge for the present
few-facet, clipping, cofactor, full-projection, or support-two results has
been performed. Reusing an earlier compiled theorem does not formalize
these new traditional arguments. No external independent kernel-checker
implementation is claimed.

The remaining unrestricted target is a dimension-free constant $C$ such
that $E(K,S)\le Cn\rho$ follows from the full actual projection information
for every prescribed maximum $S$, or a realized actual-body family proving
superlinear dimension growth. The core provides a structured comparison
matrix $(1-t)I+(t/n)J$, but its simplex need not be maximum. Face-depth
inequalities are necessary and are not sufficient for projection closeness.
Support on three or more coordinates does not reduce to the unicyclic
argument. The golden-ratio example obstructs small global constants but
is not a $c\downarrow1$ counterexample. No arbitrary determinant model is
identified with the actual geometric defect without a separate calculation.

\begin{thebibliography}{9}
\bibitem{Schneider}
R.~Schneider, \emph{Convex Bodies: The Brunn--Minkowski Theory}, second
expanded edition, Encyclopedia of Mathematics and its Applications 151,
Cambridge University Press, 2014.
\url{https://doi.org/10.1017/CBO9781139003858}.
\bibitem{Series}
Entry 005 research project, \emph{A dimension-stability structural series},
8 October 2026. Unpublished source notes and version-bound traditional
reviews. Mathematical notes, certificate scripts, provenance, and exact
source hashes are identified in the accompanying source map. No priority
certification is asserted.
\bibitem{Quadratic}
Entry 005 research project, \emph{Quadratic-dimensional simplex stability},
7 October 2026. Retained actual-body proof report and formal verification
record. The finite actual-facet specialization needed here is reproduced
in Appendices~\ref{sec:assignment}--\ref{sec:enclosure}.
\end{thebibliography}

\end{document}
